\clearpage
\fancyhead[L]{\small BEOBACHTUNGEN DER MAGNETISCHEN INKLINATION.}
\SourceLeaf{450}\pdfbookmark[0]{Beobachtungen der magnetischen Inklination in Göttingen.}{inclination-observations-title}
\begin{center}
{\Large\bfseries BEOBACHTUNGEN\par\medskip DER MAGNETISCHEN INKLINATION IN GÖTTINGEN.}
\par\bigskip
Resultate aus den Beobachtungen des magnetischen Vereins. 1841. II.
\end{center}
\begingroup
\renewcommand{\Article}[1]{\par\medskip\Needspace{4\baselineskip}\pdfbookmark[1]{#1.}{inclination-observations-art-#1}\centerline{\textbf{#1.}}\nobreak\smallskip}%
\Article{1}
Das Inklinometer, mit dem die hier mitzuteilenden Beobachtungen angestellt wurden, stammt von \textsc{Robinson}. Es war das letzte Instrument dieser Art, das der ausgezeichnete Instrumentenbauer geliefert hat.

Der vertikale Kreis hat im Lichten einen Durchmesser von $241.169$ Millimetern und ist in Abständen von zehn Minuten geteilt. Der Abstand zwischen den inneren Enden zweier Teilstriche beträgt daher $0.351$ Millimeter. Die Teilstriche erscheinen auch im Mikroskop bei beträchtlicher Vergrößerung sehr edel. Durch Messungen an mehreren von ihnen habe ich ihre Breite zu $=0.024$ Millimetern bestimmt, sodass ein Teilstrich nahezu $41$ Bogensekunden abdeckt.

Der Durchmesser des horizontalen Kreises, dort gemessen, wo das Ende des Indexstrichs die Teilstriche trifft, beträgt $148$ Millimeter. Die Teilung erfolgt in halben Graden, und der Nonius gibt einzelne Minuten an. Es wird nur eine Ablesung vorgenommen.

Die Grade des Vertikalkreises sind von beiden Endpunkten eines horizontalen Durchmessers nach oben und unten bis $90$ gezählt. Diese Einrichtung mag bei gewöhnlichen Beobachtungen vielleicht bequem erscheinen, führt aber leicht zu Verwirrung, wenn man eine absichtlich belastete Nadel verwendet und diese dadurch in einen anderen Quadranten gerät oder wenn man auch in einer zum magnetischen Meridian senkrechten Vertikalebene beobachtet. Jedenfalls macht diese Einrichtung in sol%
\SourceLeaf{451}%
chen Fällen eine etwas umständlichere und weniger übersichtliche Protokollführung nötig. Ich würde daher eine in derselben Richtung fortlaufende Graduierung von $0$ bis $360^0$ oder zweimal von $0$ bis $180^0$ vorziehen. Ich habe mir angewöhnt, im unteren Quadranten links oder im oberen rechts anstelle der eingravierten Zählung sofort die Ergänzung zu $180^0$ aufzuschreiben. In dieser Form sind alle Ablesungen in diesem Aufsatz angegeben. Auf dem Horizontalkreis laufen die Zahlen zweimal in derselben Richtung von $0$ bis $180^0$. Natürlicher und bequemer wäre eine durchgehende Zählung bis $360^0$, und in dieser Form habe ich die hier vorkommenden Ablesungen angesetzt.

An der Libelle entspricht ein Ausschlag von einem Millimeter einer Neigung von $9$ Bogensekunden.

\Article{2}
Zu dem Instrument gehören vier Nadeln, die ich durch die Zahlen $1,2,3,4$ unterscheide. Die beiden letzten haben drehbare Achsen. Auf diese Einrichtung werde ich weiter unten zurückkommen. Bei der mikroskopischen Messung war an allen acht Zapfen kein Unterschied in der Dicke zu erkennen. Ich habe diese Dicke zu $=0.590$ Millimetern bestimmt. Die Nadeln $1$ und $2$ wiegen je $16.5$ Gramm, die beiden anderen je $20.5$ Gramm.

Die einzelnen Nadeln unterscheiden sich geringfügig in ihrer Länge. Die Messung ergibt
\begin{center}
\begin{tabular}{crcrl}
für &$1$&$.\quad.\quad.\quad.\quad.\quad.$&$240.931$&Millimeter\\
&$2$&$.\quad.\quad.\quad.\quad.\quad.$&$240.866$&—\\
&$3$&$.\quad.\quad.\quad.\quad.\quad.$&$240.938$&—\\
&$4$&$.\quad.\quad.\quad.\quad.\quad.$&$240.954$&—
\end{tabular}
\end{center}
Die kürzeste Nadel ist also nur um $0.303$ und die längste nur um $0.215$ Millimeter kürzer als der lichte Durchmesser des Kreises. Das erleichtert zwar eine genauere Ablesung, hat aber zugleich zur Folge, dass schon eine sehr geringe Exzentrizität die freie Bewegung der Nadel stören kann. Daher ist es schwierig, die Teile des Instruments, von deren Stellung die Exzentrizität abhängt, ganz zufriedenstellend einzustellen, zumal ihre Stellungen noch vier weiteren, insgesamt also \textit{sechs} Bedingungen genügen müssen.

\SourceLeaf{452}\Article{3}
Diese sechs Bedingungen sind die folgenden.

Die beiden Achatplatten, auf deren oberen Kanten die Zapfen der Nadel bei der Beobachtung liegen, sollen durch die beiden Schraubenpaare, auf denen sie ruhen, so eingestellt sein, dass

1) ihre oberen Kanten in einer einzigen Ebene liegen,

2) diese Ebene senkrecht zur Ebene des Vertikalkreises steht und

3) unterhalb der Mitte dieses Kreises liegt, und zwar in einem Abstand gleich der halben Dicke eines Zapfens,

4) die Schnittlinie dieser beiden Ebenen mit der Vertikalachse einen rechten Winkel bildet.

Ferner müssen die Pfannen, mit denen man die Nadel von den Achatplatten abhebt und wieder auflegt und die mit etwas Spielraum zum Verschieben auf den Heberahmen geschraubt sind. so eingestellt sein, dass nach dem Auflegen der Nadel

5) ihre Achse senkrecht zu der zuletzt (in 4) genannten Schnittlinie steht (also zusammen mit Bedingung 2 auch senkrecht zur Ebene des Vertikalkreises) und zugleich

6) den vertikalen Durchmesser des Kreises trifft.

Die Bedingungen 1, 2, 4 zusammengenommen ersetzen eine einzige Bedingung. Bei genau senkrechter Stellung der aufrechten Drehachse soll eine horizontale Ebene die Kanten beider Achatplatten der Länge nach oder in zwei Linien berühren. Das setzt voraus, dass die Ebene des Vertikalkreises parallel zu jener Drehachse ist und daher mit ihr zugleich senkrecht steht. Man kann das als siebte Bedingung betrachten. Im Vertrauen auf die Geschicklichkeit des Instrumentenbauers setzt man sie gewöhnlich stillschweigend voraus. Das Instrument bietet in seiner jetzigen Form keine Mittel, sie zu prüfen und gegebenenfalls zu berichtigen.

\Article{4}
Bei den Beobachtungen, die in diesem Aufsatz angegeben werden, bin ich bei der Prüfung der genannten Bedingungen mangels anderer Mittel folgendermaßen vorgegangen.

\SourceLeaf{453}Zur Prüfung der \textit{ersten} Bedingung benutzte ich das Planglas eines sogenannten künstlichen Horizonts. Nachdem zuvor der Rahmen mit den Pfannen entfernt worden war, legte ich es so auf die Achatlager, dass die matt geschliffene Rückseite nach oben zeigte. Ist die Bedingung nicht erfüllt, wird stets nur \textit{eine} Achatplatte auf ihrer ganzen Länge berührt, die andere an einem Endpunkt. Das erkennt man. wenn der Fehler nicht sehr gering ist, schon mit bloßem Auge. Eine auf die Glasplatte gestellte Libelle bietet größere Genauigkeit und Sicherheit, indem sie zeigt, ob die Glasplatte zwei verschiedene Berührungslagen hat oder nur eine. Man sieht leicht, dass man mithilfe dieser Libelle nach Erfüllung der ersten Bedingung auch die \textit{zweite} und \textit{vierte} prüfen kann.

Zur Prüfung der \textit{fünften} Bedingung muss man die Nadel in zwei verschiedene Gleichgewichtslagen bringen. In beiden haben die Zapfen dieselbe Lage auf den Lagern (oder dieselbe Nadelfläche liegt vorne), und die Nadel ist nur mäßig gegen die Horizontale geneigt. Das Ende, das in der einen Lage links war, liegt aber in der anderen rechts. Am bequemsten erhält man diese beiden Lagen durch geeignete Belastungen der Nadel. Die betreffende Bedingung ist erfüllt, wenn die Schneide jedes Nadelendes in der einen Lage genauso weit vor den Kreisrand tritt wie in der anderen. In Gegenden, in denen die Inklination nur mäßig ist, ließen sich die beiden Lagen bereits durch eine bloße halbe Umdrehung des Instruments herstellen, sodass die Ebene des Kreises in beiden Fällen nahe am magnetischen Meridian liegt.

Eine ähnliche Prüfung lässt sich auch für die \textit{zweite} Bedingung anwenden. Dabei müssen allerdings zwei entgegengesetzte, nahezu vertikale Lagen der Nadel hergestellt werden. Die eine ergibt sich von selbst, wenn man die Kreisebene nahezu senkrecht zum magnetischen Meridian stellt, die andere entweder durch eine geeignete Belastung oder durch Umkehren der Pole. Man sieht aber leicht, dass dieses Verfahren mit dem oben erwähnten Verfahren mithilfe der Libelle nur dann gleichwertig ist, wenn die siebte Bedingung erfüllt ist. Durch Verbindung beider Methoden erhält man somit eine Art Prüfung dieser Bedingung selbst. Sie kann freilich nur sehr unvollkommen sein, da sich nur schätzen lässt, ob die Nadelschneide gleich weit vor den Kreisrand tritt.

\expandafter\def\csname HistoricalNoteText454\endcsname{Mit Ausnahme der vom 23. September.}%
\SourceLeaf{454}Dieselben kombinierten Lagen der Nadel dienen zugleich zur Prüfung der beiden übrigen Bedingungen. Die \textit{sechste} Bedingung ist erfüllt, wenn jedes Nadelende in der ersten nahezu horizontalen Lage genauso weit von der Innenfläche des Kreises entfernt ist wie in der zweiten. Für die \textit{dritte} Bedingung gilt Entsprechendes bei den nahezu vertikalen Lagen. Wenn die beiden Nadelhälften genau gleich lang wären, würde es zur Prüfung offenbar genügen, die Abstände beider Nadelenden von der Innenfläche des Kreises in jeweils einer nahezu horizontalen und einer nahezu vertikalen Lage miteinander zu vergleichen. Bei unserem Instrument, bei dem die Zwischenräume allgemein so klein sind, genügt das aber nicht, und selbst eine sehr geringe Ungleichheit zwischen den beiden Nadelhälften wird dabei schon bemerkbar.

\Article{5}
Wie schwierig es ist, auf diese Weise alle Bedingungen zugleich zu erfüllen, zeigt schon der Umstand, dass die beiden Schrauben, auf denen jede Achatplatte ruht, nur acht Millimeter voneinander entfernt sind. Da die Steigung eines Schraubengewindes $0.283$ Millimeter beträgt, dreht bereits eine halbe Umdrehung einer Schraube die betreffende Achatplatte um einen Grad.

Sehr erleichtert wird die Arbeit jedoch durch eine besondere Vorrichtung, die ich erst später anfertigen ließ und mit der sich die Kanten der Achatplatten in eine einzige Ebene bringen lassen, die dann waagerecht eingestellt wird. Ich halte mich jetzt aber nicht mit ihrer Beschreibung auf, da sie für die hier behandelten Beobachtungen\HistoricalNote{454} noch nicht verwendet werden konnte. Eine zweite Vorrichtung, die ebenfalls erst nach Abschluss der Beobachtungen fertig wurde, dient zur genauen Bestimmung der Abweichung des Hauptkreises von der Vertikalen. Sie ergab eine Abweichung von \textit{zehn Minuten}. Diese zu beseitigen wird jedoch erst eine Änderung am Instrument erfordern. Im Übrigen kann der Einfluss dieser Abweichung auf die Inklinationen nicht einmal eine Bogensekunde betragen.

Überhaupt muss ich darauf hinweisen, dass kleine Fehler in den verschiedenen Justierungen nur einen kaum merklichen Einfluss auf die Inklina%
\SourceLeaf{455}%
tionsbestimmungen haben können. Der Einfluss eines Teils der Fehler auf die Stellung der Nadel ist in Bezug auf diese Fehler nur von zweiter Ordnung. Die Wirkung der übrigen Fehler, insbesondere einer Exzentrizität und einer Neigung der die Achatplatten berührenden Ebene in der Richtung parallel zur Ebene des Kreises (Fehler gegen die Bedingungen 3, 6 und 4), wird durch die Kombination der einzelnen Beobachtungen vollständig beseitigt. Ich kann daher dem Urteil \textsc{Horner}'s, man müsse \textit{vor allem} auf die Beseitigung dieses letzten Fehlers achten (Physik. Wörterb. 5. Band, S. 759), nicht zustimmen. Ich halte diesen Fehler für denjenigen, auf dessen vollständige Beseitigung es am wenigsten ankommt.

\Article{6}
Die hier angegebenen Inklinationsbeobachtungen wurden alle im Freien an dem in den \textit{Resultaten} 1840. II. [S. 433 d. B.] bezeichneten Ort angestellt. Ein Schutzdach hielt die Sonnenstrahlen vom Instrument ab. Das Instrument wurde so auf dem Stein aufgestellt, dass die Gerade durch zwei Fußspitzen nahezu senkrecht zum magnetischen Meridian stand. Für diese Stellung waren die Plätze der drei Füße markiert. Die genaue magnetische Ausrichtung des Instruments wurde mithilfe einer mitgelieferten Hilfsnadel erreicht, die mit einem Achathütchen auf eine Spitze gehängt wird. Der Träger dieser Spitze hat zwei kurze zylindrische Seitenarme, die in die beiden Pfannen eingelegt werden. Durch das Gewicht des frei herabhängenden Teils des Trägers stellt sich die Spitze von selbst senkrecht. Ich habe diese Art der Ausrichtung öfter mit der sonst üblichen Methode der entsprechenden Neigungen in zwei Stellungen des Vertikalkreises kombiniert, in denen er nahezu senkrecht zum magnetischen Meridian steht. Dabei habe ich immer nur ganz unerhebliche Unterschiede gefunden. Daraus geht hervor, dass die Hilfsnadel hinreichend empfindlich ist und keine konstante Abweichung hervorruft. Eine geringe Abweichung der Vertikalebene, in der man beobachtet, vom ohnehin während der Beobachtungen nicht ganz unveränderlichen magnetischen Meridian hat im Übrigen auf die Neigung der Inklinationsnadel nur einen Einfluss zweiter Ordnung, der als ganz unmerklich gelten kann.

\SourceLeaf{456}\Article{7}
Selbst die geschicktesten Instrumentenbauer können den Schwerpunkt einer Nadel nur annähernd mit der Drehachse zur Deckung bringen. Fast immer bleibt eine Abweichung zurück. Ihren Einfluss auf die Einstellung der Nadel soll man durch Kombination von Beobachtungen unter mehrfach veränderten Umständen ermitteln oder beseitigen. Zu diesen Änderungen gehört wesentlich die Umkehrung der Nadelpole. Unter sonst gleichen Umständen ist jener Einfluss umso stärker, je schwächer die Nadel magnetisiert ist. Da man aber nicht berechtigt ist, anzunehmen. dass die Stärke des Nadelmagnetismus nach der Umkehrung der Pole wieder genauso groß ist wie vorher, hängt eine genaue Reduktion der Beobachtungen davon ab, dass man das Verhältnis dieser Stücke kennt. Man ermittelt es durch Beobachtung der Schwingungsdauer der Nadel. Ich habe jedoch aus mehreren Gründen \textit{horizontalen} Schwingungen den Vorzug gegeben und zu ihrer Beobachtung einen besonderen Apparat verwendet, den Herr Inspektor \textsc{Meyerstein} angefertigt hatte. Die Nadel schwingt in einem Holzkasten mit verglasten Deckeln. Dabei ruht sie auf einem leicht gearbeiteten Bügel, der an einem $270$ Millimeter langen, durch eine Glasröhre vor Luftzug geschützten Seidenfaden hängt. Beim Schwingen bewegen sich ihre Enden an zwei Gradbögen entlang, die jeweils $40$ Grad umfassen, in halbe Grade geteilt sind und es erlauben, $5$ Minuten sicher zu schätzen. Die Schwingungsdauer jeder Nadel wurde vor und nach dem Umstreichen jeweils aus $150$ Schwingungen bestimmt, die auf drei Reihen verteilt waren. Nach gehöriger Reduktion auf unendlich kleine Bögen ergeben diese Reihen stets vorzüglich übereinstimmende Resultate. Gewöhnlich begann man mit einem Schwingungsbogen von etwa $36$ Grad. Es verdient hier Beachtung, dass der Schwingungsbogen im Gegensatz zu den in den \textit{Resultaten} 1837. IV [S. 385 d. B.] erwähnten Erfahrungen mit schwereren Stäben an allen Tagen und bei allen Nadeln nahezu gleich schnell abnahm. Die Zeit, in der der Bogen auf seinen vierten Teil zurückging, betrug daher mit geringen Schwankungen $14$ Minuten. Übrigens wurden diese Schwingungsbeobachtungen stets in der Sternwarte auf einem Steinpostament angestellt. Es kommt dabei weniger auf die absolute Dauer an als auf das Verhältnis, das durch die geringen, möglicherweise an diesem Ort auftretenden fremden Einflüsse nicht merklich beeinflusst werden kann.

\expandafter\def\csname HistoricalNoteText457\endcsname{Das auffallendste Beispiel dieser Art wird im \textit{Fifth Report of the British association for the advancement of Science} auf S. 142 angeführt. Die acht Nadeln, mit denen Kapitän Ross in London die Inklination bestimmte, ergaben Unterschiede von bis zu 41 Minuten, obwohl die Beobachtungen mit jeder einzelnen Nadel zahlreich waren und untereinander gut übereinstimmten. Die Ursache dieser sonderbaren Erscheinung, über die keine näheren Einzelheiten mitgeteilt werden, hat man in England der nicht vollkommen zylindrischen Gestalt der Zapfen zugeschrieben. Gerade deshalb versuchte man es mit drehbaren Zapfen.}%
\SourceLeaf{457}\Article{8}
Mit den im Sommer 1842 angestellten Inklinationsbeobachtungen be\-zweckte ich neben der Feststellung der für diese Zeit geltenden magnetischen Inklination auch, den Grad der Genauigkeit zu bestimmen, den man mit dem verwendeten Instrument erreicht. Es schien mir nicht zu genügen, die Zuverlässigkeit der Endresultate, auf die so viele verschiedene Umstände Einfluss haben, anhand der Unterschiede abzuschätzen, die sich in den Einstellungen der Nadel bei wiederholtem Abheben \textit{mit dem Pfannenhebel} ergeben. Ebenso wenig kann dazu der bloße Vergleich der Inklinationsresultate aus Beobachtungen an verschiedenen Tagen dienen, da sich dabei die zufälligen, dem Instrument zuzuschreibenden Beobachtungsfehler mit den wirklichen Schwankungen der Inklination selbst vermischen. Ich war außerdem gespannt, ob meine vier Nadeln übereinstimmende Ergebnisse liefern würden oder ob sie, wie es einigen Beobachtern ergangen ist\HistoricalNote{457}, deutlich und erheblich voneinander abweichen würden.

Diese Überlegungen bewogen mich, die Beobachtungen etwas anders als üblich anzuordnen. Das Wesentliche des Unterschieds ergibt sich aus dem Folgenden.

Gewöhnlich beobachtet man den Stand der Nadel, das heißt die Stellung beider Spitzen gegenüber der Teilung des Kreises, in vier verschiedenen Kombinationen aus der Stellung des Kreises und der Art, wie die Nadel eingelegt ist. Dabei können die geteilte Fläche des Kreises und die markierte Fläche der Nadel nach Osten oder Westen, in dieselbe oder in entgegengesetzte Himmelsrichtungen weisen. Nach der Umkehrung der Pole wiederholt man dieselben Kombinationen, sodass insgesamt $16$ Ablesungen vorliegen. Sofern sich diese infolge einer starken Abweichung des Nadelschwerpunkts von ihrer Zapfenachse nicht erheblich unterscheiden, nimmt man ihr einfaches arithmetisches Mittel als Inklination. Im entgegengesetzten Fall wendet man ein aufwendigeres Rechenverfahren an. Es versteht sich, dass jede der $16$ Zahlen
\SourceLeaf{458}selbst schon das Mittel aus einer kleineren oder größeren Anzahl von Einstellungen sein kann, die man in jeder Kombination durch wiederholtes Abheben erhält.

Das von mir befolgte Verfahren unterscheidet sich davon dadurch, dass ich an jedem Tag mit zwei Nadeln beobachtete, ohne zwischen den Beobachtungen die Pole umzukehren. Die Umkehrung der Pole geschah zwischen zwei aufeinanderfolgenden Beobachtungen, und zwar abwechselnd immer nur an einer Nadel. Man sieht, dass auf diese Weise die Beobachtungen von vier Tagen alle Kombinationen der verschiedenen Polarisierungen beider Nadeln umfassen, wie es bei den Nadeln 1 und 3 vom 6. bis 9. Juli und bei den Nadeln 2 und 4 vom 17. bis 20. Juli der Fall war. Setzte man diese ähnlich kombinierten Wechsel über acht Beobachtungstage fort, wie es mit den Nadeln 1 und 2 vom 20. Mai bis 5 Juni und mit den Nadeln 3 und 4 vom 8. bis 25. Juni geschah, so erhielt man also jede Kombination der Polarisierungen zweimal. Die Beobachtungen wurden an jedem Tag so angeordnet, dass die Ergebnisse beider Nadeln möglichst gleichzeitig gewonnen wurden. Dazu beobachtete man zuerst die vier oben erwähnten Kombinationen an der einen Nadel vollständig, jede mit viermal wiederholtem Auflegen. Dann folgten die entsprechenden Kombinationen an der zweiten Nadel mit achtmal wiederholtem Auflegen. Schließlich beobachtete man an der ersten Nadel erneut dieselben Kombinationen, aber in umgekehrter Reihenfolge und mit viermal wiederholtem Auflegen.

Bei dieser Anordnung liefern die Beobachtungen eines Tages für sich allein noch keine Inklinationsbestimmung. Verbindet man sie jedoch mit den Beobachtungen des folgenden Tages, so zeigt die nicht umgestrichene Nadel offenbar, um wie viel die Inklination an den beiden Tagen verschieden war. Die einseitigen Beobachtungen an der anderen Nadel können danach auf einen einzigen Zeitpunkt reduziert und dadurch vervollständigt werden. Für eine strengere Behandlung, die sämtliche Beobachtungen aller 24 Tage umfasst, muss jedoch zuerst das Verhältnis der Teilergebnisse zueinander näher erörtert werden.

\Article{9}
Diese in mehr als einer Hinsicht wichtige Entwicklung lässt sich am bequemsten an ein Beispiel aus einer auf ge%
\SourceLeaf{459}%
wöhnliche Weise angestellten Beobachtung anknüpfen. Solche Beobachtungen habe ich auch an mehreren Tagen gemacht.

Ich wähle dazu die Beobachtung mit der Nadel 1 vom 23. September 1842 vormittags von $8\frac{1}{2}$ bis $11$ Uhr. Die magnetische Ausrichtung wurde auf die in Artikel 6 angegebene Weise mit der Hilfsnadel erreicht. Der Index des Azimutalkreises, dessen von links nach rechts zunehmende Grade ich, wie schon oben bemerkt, von $0$ bis $360^0$ durchzähle, zeigte bei der Stellung des Vertikalkreises im Meridian mit der geteilten Seite nach Osten $90^0\ 5'$.

Außer den üblichen acht Kombinationen im magnetischen Meridian beobachtete ich an diesem Tag noch ebenso viele in der zu ihm senkrechten Vertikalebene. Ich nehme diese Beobachtungen hier mit auf, da sie Gelegenheit zu mehreren Erörterungen geben. Die Nadel ist wie die drei anderen auf einer Seite an den Enden mit den Buchstaben $A$, $B$ markiert, sodass man ihre Polarisierung und die Art ihrer Einlage leicht unterscheiden kann. In jeder der $16$ Kombinationen wurde die Nadel fünfmal mit dem Pfannenhebel auf die Achatplatten gelegt. In der folgenden Übersicht gebe ich aber nur die Mittelwerte der zusammengehörigen Einstellungen an.

\begin{center}
\textit{Nadelende $B$ Nordpol.}\par\smallskip
{\fontsize{9}{11}\selectfont
\begin{tabular}{c||cc||cc}
Azimutal&\multicolumn{4}{c}{Markierte Nadelfläche}\\
Kreis&\multicolumn{2}{c||}{vorne}&\multicolumn{2}{c}{hinten}\\
&unten&oben&unten&oben\\\hline
$90^0\ 5'$&$67^0\ 27'\ 54''$&$67^0\ 29'\ 36''$&$67^0\ 45'\ 39''$&$67^0\ 44'\ 51''$\\
$180\ 5$&$89\ 52\ 39$&$89\ 52\ 51$&$90\ 12\ 30$&$90\ 10\ 30$\\
$270\ 5$&$112\ 18\ 39$&$112\ 16\ 45$&$112\ 38\ 51$&$112.\ 33\ 54$\\
$0\ 5$&$89\ 58\ 33$&$89\ 57\ 48$&$90\ 13\ 27$&$90\ 10\ 54$
\end{tabular}}
\par\medskip
\textit{Nadelende $A$ Nordpol.}\par\smallskip
{\fontsize{9}{11}\selectfont
\begin{tabular}{c||cc||cc}
$90^0\ 5'$&$68^0\ 2'\ 51''$&$68^0\ 2'\ 33''$&$67^0\ 35'\ 15''$&$67^0\ 37'\ 0''$\\
$180\ 5$&$90\ 14\ 48$&$90\ 12\ 21$&$89\ 51\ 12$&$89\ 51\ 36$\\
$270\ 5$&$112\ 27\ 21$&$112\ 22\ 33$&$112\ 7\ 6$&$112\ 5\ 33$\\
$0\ 5$&$90\ 16\ 15$&$90\ 14\ 0$&$89\ 53\ 54$&$89\ 54\ 18$
\end{tabular}}
\end{center}

\Needspace{6\baselineskip}
\SourceLeaf{460}Für die Dauer einer horizontalen Schwingung ergab sich
\begin{center}
\begin{tabular}{lll}
vor den Beobachtungen&$.\quad.\quad.\quad.\quad.\quad.\quad.\quad.\quad.$&$5''83555$\\
nach den Beobachtungen&$.\quad.\quad.\quad.\quad.\quad.\quad.\quad.\quad.$&$5.87416$
\end{tabular}
\end{center}

\Article{10}
Ich gehe nun zunächst auf die Unterschiede zwischen den Ablesungen an der unteren und der oberen Spitze ein. Sie hängen davon ab, dass die Zapfenachse weder durch den Mittelpunkt der Teilung noch durch die Gerade geht, die die beiden Spitzen der Nadel verbindet. Bezeichnen wir mit $x$, $y$ die Koordinaten des Schnittpunkts der Zapfenachse mit der Kreisebene, bezogen auf den Mittelpunkt der Teilung und ausgedrückt in Bogenteilen des inneren Kreisumfangs. Dabei sei $x$ parallel zum Durchmesser durch die beiden Nullpunkte und nach rechts positiv, $y$ parallel zum Durchmesser durch die beiden $90^0$-Punkte und nach oben positiv. Ferner bezeichne $180^0-z$ den Winkel zwischen den beiden Ebenen, die durch die Zapfenachse und die Spitzen $A$ und $B$ gelegt werden. Denkt man sich die Nadel horizontal und die markierte Seite nach oben gerichtet, so wird der Winkel von $A$ nach $B$ von links nach rechts gezählt. Schließlich sei $l$ das Mittel der beiden Ablesungen. Ihr Unterschied, verstanden als obere Ablesung minus untere Ablesung, ist dann
\[
=2x\sin l+2y\cos l\pm z
\]
Dabei gilt das obere Zeichen, wenn zugleich die markierte Fläche vorne und $A$ oben ist (hier also der Südpol) oder wenn jene Fläche hinten und $B$ oben ist. In den beiden anderen Fällen gilt das untere Zeichen.

Die obigen Beobachtungen ergeben so $16$ Gleichungen, aus denen man nach der Methode der kleinsten Quadrate erhält
\[
\begin{aligned}
x&=-38''3\\
y&=+153.2\\
z&=+75.4
\end{aligned}
\]
\Needspace{20\baselineskip}
Der Vergleich ergibt dann bei Anordnung nach der Größe von $l$
\SourceLeaf{461}\begin{center}
{\fontsize{9}{11}\selectfont
\begin{tabular}{c|r|r|r}
$l$&Beobachtung&Rechnung&Fehler\\\hline
$67^0\ 28'\ 45''$&$+102''$&$+122''$&$-20''$\\
$67\ 36\ 7$&$+105$&$+121$&$-16$\\
$67\ 45\ 15$&$-48$&$-30$&$-18$\\
$68\ 2\ 42$&$-18$&$-32$&$+14$\\
$89\ 51\ 24$&$+24$&$0$&$+24$\\
$89\ 52\ 45$&$+12$&$-1$&$+13$\\
$89\ 54\ 6$&$+24$&$-1$&$+25$\\
$89\ 58\ 10$&$-45$&$-1$&$-44$\\
$90\ 11\ 30$&$-120$&$-153$&$+33$\\
$90\ 12\ 10$&$-153$&$-153$&$0$\\
$90\ 13\ 34$&$-147$&$-153$&$+6$\\
$90\ 15\ 7$&$-135$&$-153$&$+18$\\
$112\ 6\ 19$&$-93$&$-111$&$+18$\\
$112\ 17\ 42$&$-114$&$-112$&$-2$\\
$112\ 24\ 57$&$-288$&$-264$&$-24$\\
$112\ 36\ 22$&$-297$&$-264$&$-33$
\end{tabular}}
\end{center}

Die Summe der Quadrate der verbleibenden Fehler ist $7924$. Daraus schließt man, dass man den mittleren Fehler der Differenz zweier Mittel aus je fünf Ablesungen als
\[
=\BareRoot{\frac{7924}{13}}=24''7
\]
und den mittleren Fehler einer einzelnen Ablesung an einer Spitze als
\[
=\BareRoot{\frac{5}{2}.\frac{7924}{13}}=39''0
\]
annehmen kann. Das ist in der Tat eine sehr befriedigende Genauigkeit, die eine ähnliche Auswertung der Beobachtungen von anderen Tagen nicht nur bestätigt, sondern bisweilen sogar übertrifft. Dabei sei jedoch bemerkt, dass eine solche Übereinstimmung wesentlich davon abhängt, dass man die Nadel immer nur dann abhebt, wenn sie in Ruhe oder ihrer Ruhestellung nahe ist. Ohne diese Vorsicht würde die Nadel, deren Schwingung darin besteht, dass die Zapfen auf dem Lager rollen, an einer anderen Stelle des Lagers zur Ruhe kommen als dort, wo sie aufgelegt wird. Das Exzentrizitätselement $x$ wäre also veränderlich.

Auf die hier beschriebene Weise erhält man die Werte der Exzentrizitätselemente $x$ und $y$ zwar mit großer Genauigkeit, aber diese Werte
\SourceLeaf{462}können uns nicht ohne Weiteres zeigen, ob und um wie viel die Lager und Pfannen noch verschoben werden müssen, um die Bedingungen 3 und 6 in Artikel 3 zu erfüllen. Die Bedingungen beziehen sich nämlich auf den Mittelpunkt des inneren Kreises, die Werte dagegen auf den Mittelpunkt der Teilung. Beide können etwas voneinander abweichen und tun dies bei dem betreffenden Instrument tatsächlich. Vor den hier angeführten Beobachtungen waren die betreffenden Justierungen in der Tat mit aller möglichen Sorgfalt ausgeführt worden.

Die mit $z$ bezeichnete Größe ist offenbar für jede Nadel unveränderlich, und eine ähnliche Auswertung der Beobachtungen von anderen Tagen hat nahezu denselben Wert ergeben Für die drei anderen Nadeln habe ich gefunden
\begin{center}
\begin{tabular}{crll}
für Nadel&$2$&$.\quad.\quad.\quad.\quad.\quad.\quad.\quad.$&$+3'\ 18''$\\
&$3$&$.\quad.\quad.\quad.\quad.\quad.\quad.\quad.$&$-1\ 4$\\
&$4$&$.\quad.\quad.\quad.\quad.\quad.\quad.\quad.$&$+1\ 2$
\end{tabular}
\end{center}

Obwohl die Kenntnis dieser Werte kein besonderes praktisches Interesse hat, zeugt ihre Kleinheit doch in rühmlicher Weise von der Sorgfalt, mit der der ausgezeichnete Instrumentenbauer die Nadeln bearbeitet hat.

\Article{11}
Das Mittel der Ablesungen an den beiden Spitzen gibt uns die Neigung der Geraden, die diese Spitzen verbindet, oder einer Parallelen dazu gegenüber dem mit $0$ bezeichneten Durchmesser des Vertikalkreises. Ich stelle diese $16$ Mittelwerte hier paarweise zusammen.
\begin{center}
\textit{Nadelende $B$ Nordpol}\par\smallskip
{\fontsize{9}{11}\selectfont
\begin{tabular}{c|c||c|c}
Azim.&Bez. Nadelfl.&Azim.&Bez. Nadelfl.\\
Kr.&vorne&Kr.&hinten\\\hline
$90^0\ 5'$&$67^0\ 28'\ 45''$&$270^0\ 5'$&$112^0\ 36'\ 23''$\\
$180\ 5$&$89\ 52\ 45$&$0\ 5$&$90\ 12\ 11$\\
$270\ 5$&$112\ 17\ 43$&$90\ 5$&$67\ 45\ 15$\\
$0\ 5$&$89\ 58\ 10$&$180\ 5$&$90\ 11\ 30$
\end{tabular}}
\par\medskip
\textit{Nadelende $A$ Nordpol}\par\smallskip
{\fontsize{9}{11}\selectfont
\begin{tabular}{c|c||c|c}
$90^0\ 5'$&$68^0\ 2'\ 42''$&$270^0\ 5'$&$112^0\ 6'\ 20''$\\
$180\ 5$&$90\ 13\ 34$&$0\ 5$&$89\ 54\ 6$\\
$270\ 5$&$112\ 24\ 57$&$90\ 5$&$67\ 36\ 7$\\
$0\ 5$&$90\ 15\ 8$&$180\ 5$&$89\ 51\ 24$
\end{tabular}}
\end{center}

\SourceLeaf{463}Nebeneinander stehen hier die Einstellungen, bei denen die Zapfenachse bei entgegengesetzter Lage des Vertikalkreises dieselbe Ausrichtung gegenüber den Himmelsrichtungen hatte. Der Zusammenhang zwischen zwei solchen Zahlen $l$ und $l'$ ist sehr einfach, wenn die Lager so justiert sind, dass eine zur vertikalen Drehachse senkrechte Ebene sie berührt. Unter dieser Voraussetzung liegt die Zapfenachse in beiden Einstellungen auf einer horizontalen Ebene, und die Ruhelage der Nadel ist daher offensichtlich dieselbe. Verstehen wir also unter $L$ die Neigung der Geraden von der oberen zur unteren Spitze gegenüber dem horizontalen Radius des Kreises, der jeweils rechts von der markierten Nadelfläche liegt, so hat $L$ in beiden Einstellungen denselben Wert. Diese Neigung ergibt sich aber zu
\[
\begin{aligned}
\text{in der ersten Einstellung }&=l-\alpha\\
\text{und in der zweiten }&=180^0-(l'-\alpha)
\end{aligned}
\]
wenn $\alpha$ den Fehler des Nullpunkts bezeichnet, also die Ablesung an dem Radius des Kreises, der mit der Vertikalachse einen rechten Winkel bildet. Unter der obigen Voraussetzung gilt daher
\[
\alpha=\tfrac{1}{2}(l+l')-90^0,\qquad L=\tfrac{1}{2}(l+180^0-l')
\]
Aus den Beobachtungen vom 23. September, bei denen diese Justierung mit der in Artikel 5 erwähnten Vorrichtung auf das Sorgfältigste vorgenommen worden war, erhalten wir also acht verschiedene Bestimmungen von $\alpha$
\begin{center}
\begin{tabular}{r}
$+2'\ 34''$\\
$2\ 28$\\
$1\ 29$\\
$4\ 50$\\
$4\ 31$\\
$3\ 50$\\
$0\ 32$\\
$3\ 16$
\end{tabular}
\end{center}

Die Summe der Quadrate der in Sekunden ausgedrückten Abweichungen vom Mittelwert $2'56''$ ergibt sich zu $=57214$. Betrachtet man diese Abweichungen also als rein zufällig, so liefern sie für den mittleren Fehler des
\SourceLeaf{464}Ergebnisses aus einem Paar zusammengehöriger Einstellungen $=\BareRoot{\frac{57214}{7}}=90''4$. Man sieht, dass bei diesem Instrument die Unregelmäßigkeiten beim Einstellen der Nadel viel größer sind als die reinen Ablesefehler.

\Article{12}
Anders verhält es sich jedoch, wenn die vorausgesetzte genaue Justierung der Lager fehlt. Nehmen wir an, dass ihre Kanten zwar in einer einzigen Ebene liegen, diese aber nicht senkrecht zur Vertikalachse steht. Dann ist diese Ebene in den beiden Einstellungen in entgegengesetzter Richtung gegen die Horizontalebene geneigt. Dabei kommt nur die Neigung in Richtung der Lagerkanten oder parallel zur Kreisebene in Betracht, da eine geringe Neigung in Querrichtung oder in Richtung der Nadelachse keinen merklichen Einfluss auf die Ruhelage der Nadel hat. Es bezeichne nun $L$ die Neigung der Nadel im selben Sinne wie oben, die sich beim Aufliegen auf einem vollkommen horizontalen Lager ergäbe. Mit $\delta$ bezeichnen wir die entsprechende Richtkraft, also den Koeffizienten, mit dem der Sinus einer Ablenkung von der Ruhelage multipliziert werden muss, um das Drehmoment der Kraft auszudrücken, die die Nadel in diese Stellung zurücktreibt. Schließlich sei $L+\beta$ die Neigung, die sich in der ersten Einstellung auf dem geneigten Lager tatsächlich ergibt. Dann lässt sich leicht zeigen, dass
\[
\delta\sin\beta=p\rho\sin\gamma
\]
gilt, wobei $p$ das Gewicht der Nadel, $\rho$ den Radius der Zapfen und $\gamma$ die Neigung des Lagers gegen die Horizontale bezeichnen. Die letzte Größe wird positiv genommen, wenn das Lager rechts von der markierten Nadelfläche niedriger ist. Offenbar muss man aber in der zweiten Einstellung $-\gamma$ statt $\gamma$ setzen. Dadurch geht $\beta$ in $-\beta$ über, sodass die Neigung der Nadel in der zweiten Einstellung $L-\beta$ wird. Es gilt also
\[
l-\alpha=L+\beta,\qquad180^0-(l'-\alpha)=L-\beta
\]
und damit ebenso wie im vorhergehenden Artikel
\[
\tfrac{1}{2}(l+180^0-l')=L
\]
während an die Stelle der anderen dortigen Gleichung nun
\[
\tfrac{1}{2}(l+l')-90^0=\alpha+\beta
\]


\expandafter\def\csname HistoricalNoteText465\endcsname{Genau genommen ist es die Summe der Trägheitsmomente der Nadel und des Bügels. Beide voneinander zu trennen ist teils kaum möglich, teils überflüssig, da mit diesem Bügel nur horizontale Schwingungen verwendet werden.}
\SourceLeaf{465}Liegen jedoch die Kanten der Achatplatten überhaupt nicht in einer einzigen Ebene, so bleiben diese beiden Formeln auch dann hinreichend genau gültig, wenn man nur für $\gamma$ das Mittel der Neigungen beider Kanten nimmt. Voraussetzung ist, dass der Schwerpunkt der Nadel von den beiden Auflagepunkten der Zapfen nahezu gleich weit entfernt ist. Genau genommen ergibt sich noch eine kleine Änderung daraus, dass die Gerade durch die beiden Berührungspunkte von Zapfen und Lagern in den beiden Einstellungen nicht ganz denselben Azimut hat. Der Einfluss dieses Umstands auf die Stellung der Nadel darf aber selbst dort als ganz unmerklich gelten, wo er am stärksten ist, nämlich bei Beobachtungen in der zum magnetischen Meridian senkrechten Ebene.

\Article{13}
Da es von Interesse ist zu sehen, in welchem Maße die Neigung der Lager bei fehlender Justierung auf die Einstellung der Nadel wirkt, füge ich hier noch die dafür nötigen Angaben zur am 23. September verwendeten Nadel hinzu. Um ihr Trägheitsmoment zu bestimmen, hatte ich schon früher ihre horizontalen Schwingungen beobachtet, teils ohne, teils mit einem aufgelegten Ring. Dessen eigenes Trägheitsmoment ließ sich aus Gewicht und Abmessungen mit hinreichender Genauigkeit berechnen. Am 21. September ergaben sich
\begin{center}
{\fontsize{9}{11}\selectfont
\begin{tabular}{llrl}
Schwingungsdauer ohne Ring&$.\quad.\quad.\quad.$&$5''88431$&\\
\hspace{2em}---\quad mit Ring&$.\quad.\quad.\quad.$&$7.32835$&\\
Gewicht des Ringes&$.\quad.\quad.\quad.\quad.\quad.\quad.$&$19.2385$&Gramm\\
Innerer Durchmesser&$.\quad.\quad.\quad.\quad.\quad.$&$75.525$&Millimeter\\
Äußerer Durchmesser&$.\quad.\quad.\quad.\quad.\quad.$&$79.767$&---
\end{tabular}}
\end{center}

\Needspace{6\baselineskip}
Daraus folgt, wenn man Gramm und Millimeter als Einheiten annimmt,
\begin{center}
\begin{tabular}{llr}
Trägheitsmoment des Ringes&$.\quad.\quad.\quad.$&$29019$\\
\hspace{2em}---\quad der Nadel\footnotemark&$.\quad.\quad.\quad.$&$52662$
\end{tabular}
\end{center}
\footnotetext{\csname HistoricalNoteText465\endcsname}


\SourceLeaf{466}Verbindet man dies mit den oben in Artikel 9 angegebenen Schwingungszeiten vom 23. September und setzt man die Länge des einfachen Sekundenpendels in Göttingen mit $994.126$ Millimetern an, so erhält man auf bekannte Weise
\begin{center}
\begin{tabular}{llr}
\multicolumn{3}{l}{horizontale magnetische Richtkraft}\\
\quad vor dem Umstreichen&$.\quad.\quad.\quad.$&$1.5556$\\
\quad nach dem Umstreichen&$.\quad.\quad.\quad.$&$1.5352$
\end{tabular}
\end{center}

Diese Zahlen gelten, genau genommen. zunächst nur für den Ort, an dem die Schwingungen beobachtet wurden, und schließen daher die dort möglicherweise auftretenden örtlichen Einflüsse ein. Für den gegenwärtigen Zweck spielt dieser jedenfalls nur geringe Einfluss keine Rolle.

Bei einer Neigung von $67^0\ 40'\ 54''$ ergibt sich daraus ferner
\begin{center}
\begin{tabular}{llr}
\multicolumn{3}{l}{gesamte magnetische Richtkraft,}\\
\quad vor dem Umstreichen&$.\quad.\quad.\quad.$&$4.0965$\\
\quad nach dem Umstreichen&$.\quad.\quad.\quad.$&$4.0429$\\
\multicolumn{3}{l}{vertikale magnetische Richtkraft}\\
\quad vor dem Umstreichen&$.\quad.\quad.\quad.$&$3.7897$\\
\quad nach dem Umstreichen&$.\quad.\quad.\quad.$&$3.7401$
\end{tabular}
\end{center}

Lässt man die kleine Änderung der magnetischen Richtkraft der Nadel durch die Exzentrizität ihres Schwerpunkts außer Betracht, so können diese vier Zahlen als Werte von $\delta$ gelten, je nachdem, ob im magnetischen Meridian oder in der dazu senkrechten Ebene beobachtet wurde. Da $\beta$ und $\gamma$ stets klein genug sind, um diese Größen selbst an die Stelle ihrer Sinus zu setzen, also
\[
\beta=\frac{p\rho}{\delta}.\gamma
\]
ergibt sich daraus, je nachdem, ob man die Stärke der Magnetisierung vor oder nach dem Umstreichen zugrunde legt,

für Beobachtungen im magnetischen Meridian
\[
\beta=1{,}1882\gamma\quad\text{oder}\quad\beta=1{,}2039\gamma
\]
für Beobachtungen in der zum magnetischen Meridian senkrechten Ebene
\[
\beta=1{,}2844\gamma\quad\text{oder}\quad\beta=1{,}3014\gamma
\]


\SourceLeaf{467}Die bisher betrachteten Beziehungen zwischen den einzelnen Beobachtungen sind zwar nicht wesentlich, wenn es allein darum geht, aus ihnen allen die magnetische Inklination abzuleiten. Für die Prüfung und Bestätigung des Ergebnisses sind sie jedoch nicht unwichtig, denn das rechte Vertrauen in das Ganze kann erst aus der klaren Einsicht in die befriedigende Übereinstimmung der Teile erwachsen.

\Article{14}
Die Beobachtungen sind nun auf die acht Werte von $L$ zurückgeführt. Diese lassen sich als die Neigungen der von der Südpolspitze der Nadel zur Nordpolspitze gezogenen Geraden gegenüber dem horizontalen Kreisradius rechts von der markierten Nadelfläche im Gleichgewicht verstehen. Dabei denkt man sich die Nadelzapfen auf einer horizontalen Fläche aufliegend oder nimmt an, dass sich die Nadel nur um die Achsenlinie der Zapfen drehen kann, was in statischer Hinsicht offenbar ganz dasselbe ist. Mit anderen Worten, die Werte von $L$ sind die verbesserten Werte der in Artikel 11 unter der Überschrift \textit{Bezeichnete Nadelfläche vorne} aufgeführten Zahlen. Sie sind vom Einfluss des Nullpunktfehlers und der Abweichung der Lager von der Horizontalen befreit
\begin{center}
\textit{Werte von $L$.}\par\smallskip
\begin{tabular}{c|c|c}
Az. Kr.&$B$ Nordpol&$A$ Nordpol\\\hline
$90^0\ 5'$&$67^0\ 26'\ 11''$&$67^0\ 58'\ 11''$\\
$180\ 5$&$89\ 50\ 17$&$90\ 9\ 44$\\
$270\ 5$&$112\ 16\ 14$&$112\ 24\ 25$\\
$0\ 5$&$89\ 53\ 20$&$90\ 11\ 52$
\end{tabular}
\end{center}

Um nun in einer Gleichung auszudrücken, wie die Werte von $L$ von den sie bestimmenden Größen abhängen, verwende ich die folgenden Bezeichnungen.

$V$\quad Stellung des Azimutalkreises bei der Beobachtung.

$V^0$\quad Stellung des Azimutalkreises, bei der der Vertikalkreis im magnetischen Meridian steht und die geteilte Seite nach Osten zeigt.

$i$\quad Magnetische Inklination.


\SourceLeaf{468}$m$\quad Produkt aus dem magnetischen Moment der Nadel und der gesamten Intensität der erdmagnetischen Kraft, wobei die Schwere als Einheit der beschleunigenden Kräfte angenommen wird.

$q$\quad Gewicht der Nadel, multipliziert mit dem Abstand des Schwerpunkts von der Zapfenachse.

$c$\quad Spitzer Winkel zwischen der Geraden durch die Spitzen der Nadel und ihrer magnetischen Achse, positiv, wenn diese Achse rechts liegt, sofern man sich die Nadel horizontal und mit der markierten Seite nach oben vorstellt.

$Q$\quad Winkel zwischen der Geraden von der Südpolspitze der Nadel zur Nordpolspitze und der Geraden von der Zapfenachse zum Schwerpunkt. Bei derselben Lage der Nadel wie für $c$ zählt man von der ersten Geraden aus von links nach rechts.

$\delta$\quad Richtkraft.

Zerlegt man die erdmagnetische Kraft in einen vertikalen und einen horizontalen Anteil, so ergibt sich aus dem ersten das Drehmoment, positiv in Richtung wachsender $L$ gerechnet,
\[
m\sin i\cos(L+c)
\]
aus dem anderen
\[
-m\cos i\cos(V-V^0)\sin(L+c)
\]
Die Schwere dagegen bewirkt das Drehmoment
\[
q\cos(L+Q)
\]
Da $L$ die Gleichgewichtslage bezeichnet, ist die Summe dieser drei Momente $=0$. Daraus erhalten wir die Hauptgleichung
\[
-\sin i\cos(L+c)+\cos i\cos(V-V^0)\sin(L+c)=\frac{q}{m}.\cos(L+Q)
\]
Schreiben wir in der Summe der drei Momente $L+z$ statt $L$, so erhalten wir das Drehmoment bei einer Ablenkung $z$ aus der Gleichgewichtslage. Entwickelt man diesen Ausdruck in zwei Teile mit den Faktoren $\cos z$ und $\sin z$, so verschwindet der erste aufgrund der Hauptgleichung, und der zweite wird nach dem Begriff der Richtkraft $=-\delta\sin z$.
\SourceLeaf{469}Für $\delta$ gilt also die allgemeine Formel
\[
\delta=m\sin i\sin(L+c)+m\cos i\cos(V-V^0)\cos(L+c)+q\sin(L+Q)
\]
Für die drei besonderen Hauptfälle ergibt sich daraus Folgendes.

I.\quad Für $V=V^0$
\[
\begin{aligned}
\sin(L+c-i)&=\frac{q}{m}\cos(L+Q)\\
\delta&=m\cos(L+c-i)+q\sin(L+Q)\\
&=\frac{m\cos(Q+i-c)}{\cos(L+Q)}=\frac{q\cos(Q+i-c)}{\sin(L+c-i)}
\end{aligned}
\]

II.\quad Für $V=V^0+180^0$
\[
\begin{aligned}
\sin(L+c+i)&=-\frac{q}{m}\cos(L+Q)\\
\delta&=-m\cos(L+c+i)+q\sin(L+Q)\\
&=-\frac{m\cos(Q-c-i)}{\cos(L+Q)}=\frac{q\cos(Q-c-i)}{\sin(L+c+i)}
\end{aligned}
\]

III.\quad Übereinstimmend für $V=V^0+90^0$ und $V=V^0+270^0$
\[
\begin{aligned}
\sin i\cos(L+c)&=-\frac{q}{m}\cos(L+Q)\\
\delta&=m\sin i\sin(L+c)+q\sin(L+Q)\\
&=-\frac{m\sin i\sin(Q-c)}{\cos(L+Q)}=\frac{q\sin(Q-c)}{\cos(L+c)}
\end{aligned}
\]

Unser Beispiel ergibt in den beiden letzten Fällen statt gleicher Werte von $L$ Unterschiede von $3'\ 3''$ beziehungsweise $2'\ 8''$. Sie beruhen teils auf zufälligen Beobachtungsfehlern, teils auf dem Zusammenwirken mehrerer Umstände. Die anfängliche magnetische Ausrichtung ist etwas unsicher. Die magnetische Deklination und damit der Wert von $V^0$ ändern sich im Laufe der Beobachtungen. Der Horizontalkreis hat eine geringe Exzentrizität, die sich mangels doppelter Ablesung nicht überprüfen lässt. Schließlich kann man durch Auflegen mithilfe der Pfannen nur unvollkommen erreichen, dass die Zapfenachse senkrecht zur Kreisebene steht. Alle diese Umstände macht man so weit wie möglich unschädlich, indem man die Mittelwerte aus beiden Einstellungen bildet und also


\SourceLeaf{470}\begin{center}
\begin{tabular}{llr}
für $B$ Nordpol&$.\quad.\quad.\quad.\quad.$&$L=89^0\ 51'\ 49''$\\
für $A$ Nordpol&$.\quad.\quad.\quad.\quad.$&$L=90\ 10\ 48$
\end{tabular}
\end{center}
setzt. Wegen dieser Umstände wird man dem Ergebnis für die Einstellung in einer zum magnetischen Meridian senkrechten Lage stets etwas weniger Zuverlässigkeit beimessen müssen als den Ergebnissen für die Lagen im Meridian selbst, in denen man den Einfluss jener Ursachen als unmerklich ansehen kann.

\Article{15}
Die sechs Zahlen, die uns von den ursprünglich $32$ geblieben sind, seien fortan wie folgt bezeichnet.
\begin{center}
\begin{tabular}{l|l}
Werte von $L$&für $V-V^0=$\\\hline
$f,\ f'$&$0$\\
$180^0-g,\ 180^0-g'$&$180^0$\\
$h,\ h'$&$90^0$ und $270^0$
\end{tabular}
\end{center}
Die Zeichen ohne Strich beziehen sich auf $B$ als Nordpol, die mit Strich auf $A$ als Nordpol. Für die Stellungen im magnetischen Meridian sind $f,f',g,g'$ offenbar sämtlich die Neigungen der von der Südpolspitze der Nadel zur Nordpolspitze gezogenen Geraden unter der nordlichen Horizontallinie. Die ersten beiden gelten für die Stellung, in der die markierte Nadelfläche nach Osten weist, die beiden anderen für die entgegengesetzte. Dagegen sind $h,h'$ für die Stellungen in der zum magnetischen Meridian senkrechten Ebene die Neigungen derselben Geraden zur östlichen oder westlichen Horizontallinie, je nachdem, ob die markierte Nadelfläche nach Süden oder nach Norden weist.

Von den Größen, von denen diese sechs Werte abhängen, ist $q$ ganz konstant. $i$ muss für alle als gleich angenommen werden, soweit wir die kleinen Schwankungen, die im Laufe der Beobachtungen auftreten können, nicht berücksichtigen können. $Q,m,c$ ändern dagegen nach dem Umstreichen ihre Werte. Dabei ändert sich $Q$ genau um $180^0$, während $m$ und $c$ in keinem bestimmten Zusammenhang mit ihren früheren Werten stehen, abgesehen von Folgendem. Wenn man die Pole durch gleichförmiges Streichen mit kräf\SourceLeaf{471}tigen Streichstäben umkehrt, darf man sicher sein, dass der Unterschied und für $c$ auch die absoluten Werte nicht sehr groß sein können. Lasse ich nun die Zeichen ohne Strich $Q,m,c$ die bestimmten Werte für die Beobachtungen mit $B$ als Nordpol bezeichnen und setze für die Beobachtung mit $A$ als Nordpol $Q+180^0,m',c'$ an ihre Stelle, so gehen die allgemeinen Gleichungen des vorhergehenden Artikels in die folgenden sechs über
\[
\begin{array}{lll}
\sin(f+c-i)=\frac{q}{m}\cos(f+Q)&.\ .\ .\ .\ .\ .\ .&(1)\\[4pt]
\sin(g-c-i)=\frac{q}{m}\cos(g-Q)&.\ .\ .\ .\ .\ .\ .&(2)\\[4pt]
\sin i\cos(h+c)=-\frac{q}{m}\cos(h+Q)&\phantom{.}\ .\ .\ .\ .\ .\ .&(3)\\[4pt]
\sin(f'+c'-i)=-\frac{q}{m'}\cos(f'+Q)&\phantom{.}\ .\ .\ .\ .\ .\ .&(4)\\[4pt]
\sin(g'-c'-i)=-\frac{q}{m'}\cos(g'-Q)&\phantom{.}\ .\ .\ .\ .\ .\ .&(5)\\[4pt]
\sin i\cos(h'+c')=\frac{q}{m'}\cos(h'+Q)&.\ .\ .\ .\ .\ .\ .\ .&(6)
\end{array}
\]

\Article{16}
Theoretisch reichen diese sechs Gleichungen aus, um die sechs unbekannten Größen $c,c',\frac{q}{m},\frac{q}{m'},Q,i$ zu bestimmen. Die Lösung dieser Aufgabe mag hier Platz finden, obwohl sie keinerlei praktischen Wert hat. Der enorme Einfluss der unvermeidlichen Beobachtungsfehler auf die Endresultate macht das Verfahren ganz unbrauchbar.

Multipliziert man die Gleichungen $1,2,3$ der Reihe nach mit
\[
\sin(g+h),\quad\sin(f-h),\quad\sin(f+g)
\]
und addiert sie, so erhält man nach einigen einfachen Umformungen
\[
\sin(f+c).\sin(g+h)=\sin(g-c).\sin(h-f)
\]
Daraus lässt sich $c$ leicht bestimmen, am bequemsten mit der Formel
\[
\operatorname{tang}(c+\tfrac{1}{2}f-\tfrac{1}{2}g)=-\operatorname{tang}\tfrac{1}{2}(f+g)^2.\operatorname{cotang}(h-\tfrac{1}{2}f+\tfrac{1}{2}g)
\]
Auf ähnliche Weise erhält man aus den Gleichungen $4,5,6$
\[
\operatorname{tang}(c'+\tfrac{1}{2}f'-\tfrac{1}{2}g')=-\operatorname{tang}\tfrac{1}{2}(f'+g')^2.\operatorname{cotang}(h'-\tfrac{1}{2}f'+\tfrac{1}{2}g')
\]
\SourceLeaf{472}Die Zahlen unseres Beispiels sind
\[
\begin{array}{ll}
f=67^0\ 26'\ 11''&f'=67^0\ 58'\ 11''\\
g=67\ 43\ 46&g'=67\ 35\ 35\\
h=89\ 51\ 49&h'=90\ 10\ 48
\end{array}
\]
woraus sich nach den obigen Formeln ergibt
\[
c=+12'\ 21''\qquad c'=-14'\ 18''
\]
Diese Werte sind schon so groß, dass sie kaum noch plausibel erscheinen. Wie wenig zuverlässig sie sind, wird deutlich, wenn man den Einfluss kleiner Fehler in den ihnen zugrunde liegenden Zahlen ermittelt. Der dazu dienenden Differentialformel kann man mehrere Formen geben. Eine davon lautet
\[
\begin{aligned}
\mathrm{d}c={}&-\frac{\sin(g-c).\sin(h+c)}{\sin(h-f).\sin(f+g)}\mathrm{d}f\\
&+\frac{\sin(f+c).\sin(h+c)}{\sin(g+h).\sin(f+g)}.\mathrm{d}g
+\frac{\sin(f+c).\sin(g-c)}{\sin(h-f).\sin(h+g)}.\mathrm{d}h
\end{aligned}
\]
Für $\mathrm{d}c'$ gilt dieselbe Formel, wenn man nur $f,g,h$ durch $f',g',h'$ ersetzt. Auf unsere Rechnung angewandt, ergeben die Formeln
\[
\begin{aligned}
\mathrm{d}c&=-3{,}435\mathrm{d}f+3{,}441\mathrm{d}g+5{,}876\mathrm{d}h\\
\mathrm{d}c'&=-3{,}499\mathrm{d}f'+3{,}494\mathrm{d}g'+5{,}993\mathrm{d}h'
\end{aligned}
\]
Bedenkt man also, dass die Werte von $h$ und $h'$ selbst weniger zuverlässig sind und durchaus Fehler von einer oder ein paar Minuten enthalten können, so ist klar, dass die gefundenen Werte von $c$ und $c'$ kein Vertrauen verdienen

Der Vollständigkeit halber gebe ich hier noch an, wie sich die übrigen unbekannten Größen finden lassen.

Aus der Verbindung der Gleichungen $(1)$ und $(2)$ folgt
\[
\cos i=-\frac{q}{m}.\frac{\sin(f+g)\sin(Q-c)}{\sin(2c+f-g)}\quad.\ .\ .\ .\ .\ .\ .\quad(7)
\]
und daher unter Hinzunahme von Gleichung $(3)$
\[
\operatorname{tang}i=\frac{\sin(2c+f-g)}{\sin(f+g).\cos(h+c)}.\frac{\cos(Q+h)}{\sin(Q-c)}
\]
Auf ganz ähnliche Weise ergeben die Gleichungen $4$---$6$
\expandafter\def\csname HistoricalNoteText473\endcsname{Alle in diesem Aufsatz vorkommenden Rechnungen sind zwar mit größter Genauigkeit ausgeführt worden, doch wurden beim Druck die Bruchteile der Sekunden weggelassen. Wer also mit den gekürzten Zwischenwerten weiterrechnet, wird bisweilen etwas abweichende Ergebnisse finden.}
\SourceLeaf{473}\[
\operatorname{tang}i=\frac{\sin(2c'+f'-g')}{\sin(f'+g').\cos(h'+c')}.\frac{\cos(Q+h')}{\sin(Q-c')}
\]
Schreibt man folglich zur Abkürzung
\[
\frac{\sin(2c'+f'-g').\sin(f+g).\cos(h+c)}{\sin(2c+f-g).\sin(f'+g').\cos(h'+c')}=k
\]
so gilt
\[
\cos(Q+h).\sin(Q-c')=k\cos(Q+h').\sin(Q-c)
\]
Setzt man
\[
\begin{aligned}
\cos(h-c')-k\cos(h'-c)&=A\sin B\\
\sin(h-c')-k\sin(h'-c)&=A\cos B\\
\frac{\sin(h+c')-k\sin(h'+c)}{A}&=C
\end{aligned}
\]
so nimmt diese Gleichung die einfache Form an
\[
\cos(2Q-B)=C
\]
Dadurch ist $Q$ bestimmt. Dann findet man $i$ aus einer der beiden Gleichungen für $\operatorname{tang}i$ und schließlich $\frac{q}{m}$ und $\frac{q}{m'}$ aus $(1)$ oder $(2)$ und aus $(4)$ oder $(5)$. Zu diesen Rechnungen ist noch Folgendes zu bemerken.

I.\quad Um die Zahlenrechnung nach den obigen Formeln mit Genauigkeit ausführen zu können, muss man $c$ und $c'$ viel genauer berechnen, als es ihre absolute Unzuverlässigkeit an sich rechtfertigt. Andernfalls würden die beiden Bestimmungen von $i,\frac{q}{m},\frac{q}{m'}$ nur schlecht übereinstimmen\HistoricalNote{473}. Für jene Formeln kann man andere einsetzen, die diesen Nachteil vermeiden, aber etwas weniger einfach sind. Ich gebe sie hier an und lasse die nicht schwierige Herleitung weg.
\[
\begin{aligned}
\operatorname{tang}i&=-\frac{2\sin(f+c).\sin(g-c)}{\sin(f+g).\sin(h+c)}.\frac{\cos(Q+h)}{\sin(Q-c)}\\
&=-\frac{2\sin(f'+c').\sin(g'-c')}{\sin(f'+g').\sin(h'+c')}.\frac{\cos(Q+h')}{\sin(Q-c')}
\end{aligned}
\]
\[
k=\frac{\sin(f+g).\sin(f'+c').\sin(g'-c').\sin(h+c)}{\sin(f'+g').\sin(f+c).\sin(g-c).\sin(h'+c')}
\]

\SourceLeaf{474}II.\quad Die Gleichung $\cos(2Q-B)=C$ hat, abgesehen vom Sonderfall $C=\pm1$, stets vier verschiedene Lösungen, also zwischen $0$ und $360^0$ liegende Werte von $Q$, die sich paarweise um $180^0$ unterscheiden. Zwei solche Werte von $Q$ gehören zum selben Wert von $i$, aber zu Werten von $\frac{q}{m},\frac{q}{m'}$, die dem Betrag nach gleich sind und entgegengesetzte Vorzeichen haben. Da diese letzteren Größen ihrer Natur nach positiv sein müssen, entfällt in jedem Paar von selbst ein Wert von $Q$. Liefert aber ein Wert von $Q$ für $\frac{q}{m},\frac{q}{m'}$ untereinander entgegengesetzte Vorzeichen, so muss man offenbar das ganze Paar verwerfen. Sollte das bei beiden Paaren der Fall sein, so lässt sich daraus nur schließen, dass die Beobachtungsfehler die Verbindung der Gleichungen $1$---$6$ zur Bestimmung der unbekannten Größen ganz untauglich machen. In unserem Beispiel ergibt die Rechnung folgende. zwei Systeme von Werten
\begin{center}Erstes System\end{center}
\[
\begin{aligned}
Q&=\left\{\begin{array}{rrr}12^0&44'&41''\\192&44&41\end{array}\right.\\
i&=\phantom{\{}\begin{array}{rrr}67&41&33\end{array}\\
\frac{q}{m}&=\mp0.0051395\\
\frac{q}{m'}&=\mp0.0042073
\end{aligned}
\]
\begin{center}Zweites System\end{center}
\[
\begin{aligned}
Q&=\left\{\begin{array}{rrr}179^0&57'&42''\\359&57&42\end{array}\right.\\
i&=\phantom{\{}\begin{array}{rrr}60&2&11\end{array}\\
\frac{q}{m}&=\mp0.3443905\\
\frac{q}{m'}&=\pm0.3563855
\end{aligned}
\]

Hier muss man offenbar das zweite System ganz verwerfen und im ersten den oberen Wert von $Q$. Somit ist nur der Wert $Q=192^0\ 44'\ 41''$ zulässig. Dass damit ein recht guter Wert von $i$ verbunden ist und dass das Verhältnis der Werte von $\frac{q}{m}$ und $\frac{q}{m'}$ vom Verhältnis der Quadrate der Schwingungszeiten (Artikel\ 9), denen jene Werte proportional sein sollten, schon sehr stark, aber nicht noch viel stärker abweicht, ist bloß einem zufälligen Ausgleich der Beobachtungsfehler zuzuschreiben. Tatsächlich liefert bereits eine bloße Vergrößerung des Werts von $h'$ um eine Minute bei unveränderten Werten
\SourceLeaf{475}der übrigen fünf Größen $f,g,h,f',g'$ ganz untaugliche Ergebnisse. Die Rechnung nach dem obigen Verfahren ergibt zwei Lösungssysteme, in denen die Neigung $68^0\ 17'\ 40''$ beziehungsweise $66^0\ 23'\ 12''$ beträgt, während in beiden Systemen die Werte von $\frac{q}{m},\frac{q}{m'}$ entgegengesetzte Vorzeichen erhalten. Das ist ein schlagender Beweis dafür, dass man die Rechnung nicht auf solche Kombinationen gründen darf.

\Article{17}
Lassen wir nun die Beobachtungen in der zum magnetischen Meridian senkrechten Ebene beiseite, so müssen wir sie entweder durch andere Daten ersetzen oder gewisse willkürliche Annahmen zugrunde legen, die nicht streng richtig sind, und uns mit der so erreichbaren Genauigkeit der Ergebnisse begnügen. Bei meinen Beobachtungen lässt sich durchweg eine neue Angabe aus den vor und nach der Umkehrung der Pole bestimmten Schwingungsdauern gewinnen. Ihre Quadrate können als proportional zu den Größen $\frac{q}{m},\frac{q}{m'}$ betrachtet werden. Derselbe Apparat, mit dem diese Schwingungsdauern beobachtet werden, kann zwar auch zur unmittelbaren Bestimmung von $c$ und $c'$ dienen. Dazu legt man die Nadel zweimal in den Bügel, einmal mit der markierten Seite nach oben und einmal nach unten, beobachtet die Stellung der Spitzen gegenüber dem Gradbogen und berücksichtigt etwaige Änderungen der Deklination durch gleichzeitige Beobachtungen am Unifilarmagnetometer. Der Apparat erlaubt aber keine so genauen Ablesungen, wie sie für \textit{diese} Anwendung erforderlich wären (für die er nicht bestimmt ist). Wäre ein solcher Apparat viel feiner geteilt, wäre für eine unverrückbare Aufstellung gesorgt und würde etwa mit Mikroskopen abgelesen, so ließen sich $c$ und $c'$ allerdings unmittelbar mit aller wünschenswerten Genauigkeit bestimmen. Wir hätten dann sogar eine Angabe mehr als nötig, sodass eine angemessene Ausgleichung die Genauigkeit des Ergebnisses noch erhöhen könnte.

Ich ersetze daher vorerst die fehlende Angabe durch die Annahme, dass sich die magnetische Achse der Nadel durch die Umkehrung der Pole nicht verändert, oder dass $c'=c$. Diese Annahme haben alle Beobachter gemacht, die die Inklination mit einer strengeren Rechnung als nach der sonst allgemein gebräuchlichen Formel $i=\frac14(f+g+f'+g')$ zu bestimmen versucht
\SourceLeaf{476}haben. Man hat allerdings Grund anzunehmen, dass diese Annahme kaum je \textit{weit} von der Wirklichkeit abweicht, wenn man das Streichen immer mit großer Sorgfalt, mit denselben Streichstäben und bei gleicher Lage der Nadel in einem zweckmäßig gebauten Trog ausführt. Meine eigenen Erfahrungen zeigen jedoch, dass trotz dieser Vorsicht nicht unbedeutende Unterschiede in der Lage der magnetischen Achse der Nadel vorkommen können. Auch in den Angaben anderer Beobachter erkennt man oft sichere Spuren davon. (So ergeben zum Beispiel \textsc{Ermans} Beobachtungen vom 13. Oktober 1829, nach seinen eigenen Grundsätzen behandelt, für die eine Nadel eine Abweichung der magnetischen Achse von $36'\ 24''$, während sie zu anderen Zeiten sehr klein gewesen zu sein scheint). Glücklicherweise kann unter Umständen wie den hier vorliegenden selbst ein beträchtlicher Fehler in jener Annahme nur einen sehr geringen Einfluss auf das Ergebnis haben.

\Article{18}
Auf dieser Grundlage lässt sich die Aufgabe wie folgt lösen. Mit der bereits oben verwendeten Gleichung $(7)$
\[
\frac{\cos i.\sin(2c+f-g)}{\sin(f+g)}=-\frac{q}{m}.\sin(Q-c)
\]
verbinde ich die Gleichung, die sich auf ähnliche Weise aus $(4)$ und $(5)$ ergibt. Darin schreibe ich $c$ statt $c'$ und $\frac{\lambda q}{m}$ statt $\frac{q}{m'}$,
\[
\frac{\cos i.\sin(2c+f'-g')}{\sin(f'+g')}=-\frac{\lambda q}{m}.\sin(Q-c)\quad.\ .\ .\ .\ .\ .\quad(8)
\]
also
\[
\lambda\sin(f'+g')\sin(2c+f-g)=\sin(f+g).\sin(2c+f'-g')
\]
wodurch $c$ bestimmt ist, am bequemsten mit der Formel $(9)$
\[
\operatorname{tang}(2c-\tfrac12(g+g'-f-f'))=\frac{\lambda\sin(f+g)-\sin(f'+g')}{\lambda\sin(f+g)+\sin(f'+g')}.\operatorname{tang}(f-g-f'+g')
\]

Ferner folgt aus $(1)$ und $(2)$
\[
2\cos i.\sin(f+c)\sin(g-c)-\sin i.\sin(f+g)=\frac{q}{m}.\cos(Q-c)
\]
und daraus in Verbindung mit $(7)$
\[
\operatorname{cotang}(Q-c)=\frac{\sin(f+g)}{\sin(2c+f-g)}.\operatorname{tang}i-\frac{2\sin(f+c).\sin(g-c)}{\sin(2c+f-g)}
\]
\SourceLeaf{477}Auf ähnliche Weise ergibt sich aus $(4)$, $(5)$ und $(8)$
\[
\operatorname{cotang}(Q-c)=\frac{\sin(f'+g')}{\sin(2c+f'-g')}.\operatorname{tang}i-\frac{2\sin(f'+c).\sin(g'-c)}{\sin(2c+f'-g')}
\]
Schreibt man zur Abkürzung
\[
\begin{array}{ll}
\operatorname{cotang}(f+c)=F&\operatorname{cotang}(f'+c)=F'\\
\operatorname{cotang}(g-c)=G&\operatorname{cotang}(g'-c)=G'
\end{array}
\]
so nehmen diese beiden Gleichungen die Form an
\[
\begin{aligned}
\operatorname{cotang}(Q-c)&=\frac{G+F}{G-F}.\operatorname{tang}i-\frac{2}{G-F}\\
\operatorname{cotang}(Q-c)&=\frac{G'+F'}{G'-F'}.\operatorname{tang}i-\frac{2}{G'-F'}
\end{aligned}
\]
woraus sich schließlich ergibt
\[
\begin{aligned}
\operatorname{tang}i&=\frac{G'-F'-G+F}{G'F-GF'}\\
\operatorname{cotang}(Q-c)&=\frac{G'+F'-G-F}{G'F-GF'}
\end{aligned}
\]
Nachdem $i$ und $Q$ gefunden sind, kann man $\frac{m}{q}$ aus einer beliebigen der Gleichungen $1,2,4,5,7,8$ bestimmen.

In unserem Beispiel haben wir
\[
\lambda=\left(\frac{5{,}87416}{5{,}83555}\right)^2
\]
und die weitere Rechnung ergibt
\[
\begin{array}{rl}
c=&-0^0\ 1'\ 13''\\
i=&67\ 40\ 54\\
Q-c=&145\ 17\ 10\\
Q=&145\ 15\ 57\\
\dfrac{q}{m}=&0.0055111\\[4pt]
\dfrac{q}{m'}=&0.0055843
\end{array}
\]
Die aus diesen Elementen \textit{berechneten} Werte von $h,h'$ sind
\[
\begin{array}{rl}
h=&89^0\ 49'\ 30''\\
h'=&90\ 12\ 59
\end{array}
\]
von diesen weichen die beobachteten Werte also um $+2'\ 19''$ und $-2'\ 11''$ ab.

\SourceLeaf{478}\Article{19}
Wenn das Verhältnis von $\frac{q}{m},\frac{q}{m'}$ nicht unmittelbar bestimmt wurde, muss man statt einer willkürlichen Annahme zwei machen. Die Beobachter haben die folgenden beiden Gruppen von Annahmen verwendet.

I.\quad Man nimmt an, dass zugleich $c=0$ und $c'=0$ gilt. Dann haben wir für $i$ die Formel
\[
\operatorname{tang}i=\frac{\operatorname{cotg}g'-\operatorname{cotg}f'-\operatorname{cotg}g+\operatorname{cotg}f}{\operatorname{cotg}g'\operatorname{cotg}f-\operatorname{cotg}f'\operatorname{cotg}g}
\]
Dies ist das übliche Verfahren, wenn man nach \textsc{Mayers} Vorbild die Nadel absichtlich mit einem kleinen seitlichen Gewicht belastet hat. So erhält man Einstellungen der Nadel an ganz anderen Stellen des Teilkreises als ohne Belastung. Erhält man dabei keine wesentlich abweichenden Ergebnisse, so kann man einigermaßen beruhigt sein, dass der Teilkreis keine Teile mit eigenem Magnetismus enthält. Es ist allerdings ratsam, die Belastung mäßig zu halten. Bei größerer Belastung würden die Beobachtungsfehler einen übermäßig großen Einfluss auf das Ergebnis haben, und auch die Vernachlässigung von $c,c'$ würde noch eine merkliche nachteilige Wirkung haben.

II.\quad Man setzt voraus, dass $m=m'$ und $c=c'$ gilt. Das ist ersichtlich nur ein Sonderfall des im vorhergehenden Artikel behandelten Falls. Man kann die dortigen Formeln also unmittelbar anwenden, indem man $\lambda=1$ setzt. Die Formel $(9)$ für $c$ nimmt dann eine noch etwas einfachere Form an
\[
\operatorname{tang}(2c-\tfrac12(g+g'-f-f'))=\frac{\operatorname{tang}\frac12(f+g-f'-g').\operatorname{tang}\frac12(f-g-f'+g')}{\operatorname{tang}\frac12(f+g+f'+g')}
\]
Wird $c$ nicht verlangt und soll lediglich $i$ bestimmt werden, findet sich eine elegante Rechenvorschrift in \textsc{Ermans} Reise, 2.\ Abteilung, 2.\ Band, S.\ 22.

\Article{20}
Die bisher entwickelten Beziehungen der Beobachtungen zur Inklination und zu den übrigen Elementen gelten allgemein, unabhängig davon, ob die Abweichung des Schwerpunkts von der Zapfenachse groß oder klein ist. Letzteres ist bei Nadeln von einem tüchtigen Instrumentenbauer immer der Fall, solange sie nicht durch äußere Ursachen (zum Beispiel Rostflecken, Abschleifen, Entfernen der Zapfen oder absichtlich angebrachte Zusatzgewichte) verän\SourceLeaf{479}dert werden. Dann lassen sich die Formeln sehr wesentlich vereinfachen. Solange $\frac{q}{m}$ oder $\frac{q}{m'}$ den Wert $0.03$ nicht überschreitet, kann der Unterschied zwischen den Sinuswerten von $f+c-i,\ g-c-i,\ f'+c'-i,\ g'-c'-i$ und den Bögen selbst keine Sekunde erreichen. Angesichts der mäßigen Genauigkeit, die Beobachtungen mit dem Inklinometer erlauben, darf man daher selbst bei erheblich größeren Werten von $\frac{q}{m},\frac{q}{m'}$ Bogen und Sinus bedenkenlos vertauschen. Bei den vier Nadeln des \textsc{Robinson}schen Inklinometers liegen die Werte in noch viel engeren Grenzen. Daher werde ich die hier mitzuteilenden Beobachtungen nach einem solchen verkürzten Verfahren behandeln und es vorher auf das bisher betrachtete Beispiel anwenden.

\Article{21}
Setzen wir zur Abkürzung
\[
\begin{aligned}
\frac{206265''q\cos Q}{m}&=t\\
\frac{206265''q\sin Q}{m}&=u
\end{aligned}
\]
so nehmen unter der Voraussetzung, dass
\[
f+c-i,\quad g-c-i,\quad f'+c'-i,\quad g'-c'-i
\]
klein genug sind, um sie mit ihren Sinuswerten zu vertauschen, die Gleichungen $1,2,4,5$ aus Artikel 15 die folgende Form an
\[
\begin{aligned}
i&=f+c-t\cos f+u\sin f\\
i&=g-c-t\cos g-u\sin g\\
i&=f'+c'+\lambda t\cos f'-\lambda u\sin f'\\
i&=g'-c'+\lambda t\cos g'+\lambda u\sin g'
\end{aligned}
\]
Die fünf unbekannten Größen $i,c,c',t,u$ lassen sich zwar durch vier Gleichungen nicht bestimmen, wohl aber durch eine einzige unbestimmt bleibende Größe ausdrücken. Wählt man dafür $c'-c$, so wird besonders deutlich, in welchem Maß man diese Größe vernachlässigen darf. Die Elimination selbst führt man in jedem Einzelfall am bequemsten erst durch, nachdem man die Zahlenwerte der Beobachtungsdaten eingesetzt hat.

\SourceLeaf{480}In unserem Beispiel lauten die vier Gleichungen dann
\[
\begin{aligned}
i&=67^0\ 26'\ 11''+c-0{,}3837t+0{,}9234u\\
i&=67\ 43\ 46-c-0{,}3790t-0{,}9254u\\
i&=67\ 58\ 11+c'+0{,}3801t-0{,}9393u\\
i&=67\ 35\ 35-c'+0{,}3862t+0{,}9368u
\end{aligned}
\]
woraus man durch Elimination erhält
\[
\begin{aligned}
i&=67^0\ 41'\ 54''-0{,}0006(c'-c)\\
t&=-934+0{,}0002(c'-c)\\
u&=+648+0{,}5369(c'-c)\\
\tfrac12(c'+c)&=-73+0{,}0037(c'-c)
\end{aligned}
\]

Daraus erkennt man, dass die willkürliche Annahme, $c$ und $c'$ seien gleich, eine zuverlässige Bestimmung von $u$ unmöglich macht. Sie hat aber keinen merklichen Einfluss auf die Werte von $i$ und $t$ und selbst auf die Bestimmung des Mittelwerts von $c$ und $c'$ nur einen geringen.

Das Mittel aus den vier Gleichungen ist
\[
i=67^0\ 40'\ 56''+0{,}0009t-0{,}0011u
\]
wobei das absolute Glied das einfache Mittel aus $f,g,f',g'$ ist und daher ohne Weiteres als Inklination hätte angenommen werden können. Das ist tatsächlich das übliche Verfahren. Es ist auch immer unbedenklich, wenn die Werte von $f,g,f',g'$ keine großen Unterschiede aufweisen.

\Article{22}
Bevor ich das bisher behandelte Beispiel verlasse, möchte ich noch bemerken, dass sich die Gleichungen $3$ und $6$ ganz ähnlich vereinfachen lassen wie die anderen. Man kann nämlich setzen
\[
\begin{aligned}
c&=90^0-h+\frac{\cos h}{\sin i}.t-\frac{\sin h}{\sin i}.u\\
c'&=90^0-h'-\frac{\lambda\cos h'}{\sin i}.t+\frac{\lambda\sin h'}{\sin i}.u
\end{aligned}
\]
Bei der Zahlenrechnung kann man hier für $i$ bedenkenlos den Wert
\SourceLeaf{481}$\frac14(f+g+f'+g')$ einsetzen. Damit nehmen diese Gleichungen in unserem Beispiel die folgende Form an
\[
\begin{aligned}
c&=+491''+0{,}0026t-1{,}0810u\\
c'&=-648+0{,}0034t+1{,}0953u
\end{aligned}
\]
Die Werte von $h$ und $h'$ beruhen auf doppelt so vielen Einstellungen wie die Werte von $f,g,f',g'$. Käme es allein auf die Anzahl der Einstellungen an, müsste man daher jeder dieser Gleichungen das Gewicht $2\sin i^2$ geben und das Gewicht jeder der vier Gleichungen des vorhergehenden Artikels $=1$ setzen. Aus den oben (Artikel 14) angeführten Gründen sind die Bestimmungen von $h,h'$ jedoch erheblich weniger zuverlässig. Zur Vereinfachung der Rechnung kann man deshalb für alle sechs Gleichungen dasselbe Gewicht annehmen. Berechnet man aus ihnen auf diese Weise die fünf unbekannten Größen nach der Methode der kleinsten Quadrate, so erhält man
\[
\begin{array}{rl}
i=&67^0\ 40'\ 55''\\
t=&-934\\
u=&-211\\
c=&+719\\
c'=&-880
\end{array}
\]
Diese Werte erfüllen alle Gleichungen bis auf $1''$ und $2''$. Ein solches Maß an Übereinstimmung kann man freilich nur als zufällig ansehen, da die Daten sehr viel unzuverlässiger sind. Auch die Werte von $u,c,c'$ verdienen kein Vertrauen. Denn bei so großen Inklinationen wie im unsern Gegenden eignen sich die Daten überhaupt nicht dazu, diese Größen auch nur einigermaßen zuverlässig voneinander zu unterscheiden.

\Article{23}
Nach dieser Durchsicht der verschiedenen Rechenmethoden gehe ich zum Hauptgegenstand über und stelle zunächst die Beobachtungen tabellarisch zusammen, die auf die in Artikel 8 beschriebene Weise angestellt wurden. Ich führe hier nur die mit $f,g,f',g'$ bezeichneten Größen auf und lasse die Teilergebnisse weg, aus denen sie nach dem in den Artikeln $11$---$13$ angegebenen Verfahren abgeleitet wurden. Das geschieht teils aus Platzgründen, teils weil die Elemente, mit denen sie zusammenhängen,
\SourceLeaf{482}wegen häufiger Änderungen an den Lagern und Pfannen an den verschiedenen Tagen unterschiedliche Werte hatten. Meistens wurden die Beobachtungen vormittags zwischen 8 und 11 Uhr angestellt. Am 16., 22. und 25. Juni und am 17. und 20. Juli wurden sie jedoch nachmittags zwischen 4 und 6 Uhr angestellt.

Die einzelnen Spalten geben das Zeichen des Nordpolendes der Nadel an, die Werte von $f$ und $g$ oder von $f'$ und $g'$, je nachdem, ob $B$ oder $A$ der Nordpol war, sowie die Dauer der horizontalen Schwingung.

\Needspace{16\baselineskip}
\begin{center}
\textit{Beobachtungen mit Nadel} 1.

\begin{tabular}{rr|c|rrr|rrr|l}
Mai&20&$B$&$67^0$&$11'$&$0''$&$67^0$&$58'$&$46''$&$5''87152$\\
&21&$A$&&$57$&$1$&&$35$&$14$&$5.81508$\\
&22&$A$&&$56$&$29$&&$36$&$45$&$5.82044$\\
&24&$B$&&$16$&$45$&&$45$&$48$&$5.81557$\\
&31&$B$&&$18$&$1$&&$49$&$41$&$5.82075$\\
Juni&2&$A$&&$53$&$55$&&$33$&$9$&$5.85778$\\
&4&$A$&&$56$&$38$&&$32$&$10$&$5.86442$\\
&5&$B$&&$24$&$13$&&$46$&$44$&$5.83615$\\
Juli&6&$A$&&$59$&$41$&&$35$&$21$&$5.83716$\\
&7&$A$&&$58$&$7$&&$37$&$51$&$5.83818$\\
&8&$B$&&$20$&$8$&&$44$&$47$&$5.89602$\\
&9&$B$&&$20$&$43$&&$44$&$25$&$5.90035$
\end{tabular}
\end{center}

\Needspace{16\baselineskip}
\begin{center}
\textit{Beobachtungen mit Nadel} 2.

\begin{tabular}{rr|c|rrr|rrr|l}
Mai&20&$A$&$67^0$&$40'$&$57''$&$67^0$&$20'$&$37''$&$5''72416$\\
&21&$A$&&$41$&$8$&&$21$&$5$&$5.72453$\\
&22&$B$&&$43$&$28$&&$50$&$45$&$5.65355$\\
&24&$B$&&$41$&$43$&&$54$&$32$&$5.66875$\\
&31&$A$&&$43$&$34$&&$18$&$29$&$5.67439$\\
Juni&2&$A$&&$41$&$46$&&$18$&$12$&$5.67665$\\
&4&$B$&&$42$&$42$&&$46$&$57$&$5.68010$\\
&5&$B$&&$44$&$53$&&$50$&$24$&$5.68890$\\
Juli&17&$B$&&$45$&$20$&&$50$&$17$&$5.70183$\\
&18&$A$&&$40$&$26$&&$22$&$50$&$5.68692$\\
&19&$A$&&$40$&$21$&&$22$&$10$&$5.69677$\\
&20&$B$&&$40$&$40$&&$54$&$19$&$5.66585$
\end{tabular}
\end{center}

\SourceLeaf{483}\Needspace{16\baselineskip}
\begin{center}
\textit{Beobachtungen mit Nadel} 3.

\begin{tabular}{rr|c|rrr|rrr|l}
Juni&8&$B$&$67^0$&$47'$&$58''$&$67^0$&$48'$&$52''$&$6''17149$\\
&9&$B$&&$40$&$55$&&$42$&$28$&$6.18077$\\
&11&$A$&&$30$&$58$&&$32$&$35$&$6.18080$\\
&16&$B$&&$40$&$0$&&$42$&$40$&$6.17046$\\
&18&$B$&&$43$&$13$&&$47$&$40$&$6.18005$\\
&22&$A$&&$27$&$33$&&$39$&$19$&$6.16591$\\
&23&$A$&&$29$&$46$&&$41$&$8$&$6.16948$\\
&25&$A$&&$29$&$3$&&$41$&$7$&$6.17663$\\
Juli&6&$A$&&$32$&$38$&&$40$&$37$&$6.18305$\\
&7&$B$&&$45$&$56$&&$42$&$12$&$6.17982$\\
&8&$B$&&$46$&$59$&&$43$&$37$&$6.18339$\\
&9&$A$&&$30$&$42$&&$39$&$42$&$6.23905$
\end{tabular}
\end{center}

\Needspace{16\baselineskip}
\begin{center}
\textit{Beobachtungen mit Nadel} 4.

\begin{tabular}{rr|c|rrr|rrr|l}
Juni&8&$A$&$67^0$&$45'$&$9''$&$67^0$&$27'$&$3''$&$5''96200$\\
&9&$B$&&$22$&$56$&$68$&$8$&$28$&$5.91653$\\
&11&$B$&&$23$&$16$&&$7$&$48$&$5.94665$\\
&16&$A$&&$49$&$54$&$67$&$12$&$8$&$6.01785$\\
&18&$B$&&$27$&$48$&$68$&$8$&$45$&$5.93204$\\
&22&$B$&&$26$&$46$&&$3$&$56$&$5.94065$\\
&23&$A$&&$50$&$19$&$67$&$15$&$37$&$5.93939$\\
&25&$A$&&$50$&$4$&&$15$&$22$&$5.94731$\\
Juli&17&$A$&&$50$&$13$&&$15$&$43$&$5.96850$\\
&18&$A$&&$49$&$57$&&$14$&$48$&$5.96931$\\
&19&$B$&&$22$&$43$&$68$&$9$&$18$&$5.92673$\\
&20&$B$&&$22$&$41$&&$10$&$19$&$5.92783$
\end{tabular}
\end{center}

\Article{24}
Bei der Berechnung dieser Beobachtungen werde ich anstelle der oben (Artikel 21 und 22) verwendeten $t,u$ etwas abgeänderte Hilfsgrößen einführen. Bezeichnet man für eine der Nadeln die Dauer einer horizontalen Schwingung mit $n$, die Summe der Trägheitsmomente der Nadel und des Bügels bezüglich der bei diesen Schwingungen vertikalen Drehachse mit $k$ und die Länge des einfachen Sekundenpendels mit $l$. so gilt bekanntlich
\[
lmnn\cos i=k
\]
\SourceLeaf{484}Man wähle eine Normalschwingungsdauer $N$ und eine Normalinklination, die ungefähr in der Mitte der vorgekommenen Werte von $n$ beziehungsweise $i$ liegen, und bezeichne den entsprechenden Wert von $m$ mit $M$, sodass
\[
lMNN\cos I=k
\]
gilt. Schließlich sei
\[
\begin{aligned}
x&=\frac{q\cos Q.\cos I.206265''}{M}\\
y&=\frac{q\sin Q.\sin I.206265''}{M}
\end{aligned}
\]
Diese Größen sind also für alle Beobachtungen mit dieser Nadel konstant. Die Gleichungen lauten dann
\[
\begin{aligned}
i&=f+c-\frac{nn\cos f}{NN\cos I}.\frac{\cos i}{\cos I}.x+\frac{nn\sin f}{NN\sin I}.\frac{\cos i}{\cos I}.y\\
i&=g-c-\frac{nn\cos g}{NN\cos I}.\frac{\cos i}{\cos I}.x-\frac{nn\sin g}{NN\sin I}.\frac{\cos i}{\cos I}.y
\end{aligned}
\]
wenn $B$ der Nordpol ist. Ist $A$ der Nordpol, muss man nur die Vorzeichen der Glieder umkehren, die $x$ und $y$ enthalten.

Diese Form hat den Vorteil, dass die Koeffizienten von $x$ und $y$ immer nur wenig von eins abweichen. Bei einer so geringen Exzentrizität des Schwerpunkts wie bei den vier fraglichen Nadeln und bei so mäßigen Schwankungen von $n$ kann man anstelle jener Koeffizienten ohne Weiteres eins setzen. Das nenne ich die abgekürzte Rechnung. Ich habe mir dennoch die Mühe gemacht, die $192$ Koeffizienten genauer zu berechnen und nur den Faktor $\frac{\cos i}{\cos I}$ wegzulassen, auch wenn der Nutzen hauptsächlich darin besteht, die Zulässigkeit der abgekürzten Rechnung deutlicher zu zeigen. Fortan sollen sich die Buchstaben $N,x,y$ ohne Strich auf Nadel $1$ beziehen. Die Werte für die drei anderen Nadeln werden der Reihe nach durch einen, zwei und drei Striche unterschieden. Für die vorliegende Rechnung sind die folgenden Werte gewählt
\[
\begin{aligned}
I&=67^0\ 40'\ 0''\\
N&=5''847785\\
N'&=5.686867\\
N''&=6.181742\\
N'''&=5.949567
\end{aligned}
\]
\expandafter\def\csname HistoricalNoteText485\endcsname{Es versteht sich, dass ein solcher Fehler, sofern er überhaupt vorliegt, nur auf eine Abweichung der Zapfen von der Zylinderform zurückgeführt werden kann. Er bleibt nur so lange konstant, wie die Zapfen stets an denselben Stellen aufliegen. Bei einer ganz anderen Inklination könnte er also auch einen ganz anderen Wert haben.}
\SourceLeaf{485}Um Platz zu sparen, werde ich die Rechnungen selbst nur so weit mitteilen, wie es nötig ist, um ihrem Gang im Allgemeinen folgen zu können. Am stärksten von eins weichen übrigens die Koeffizienten $0.96895$ und $1.04324$ ab, die am 9. und 16. Juni bei Nadel 4 vorkommen.

\Article{25}
Die Beobachtungen mit einer Nadel liefern an jedem Tag zwei Gleichungen. Halbiert man ihre Summe und ihre Differenz, so erhält man zwei weitere Gleichungen, die mit I und II bezeichnet seien. Es entstehen also $48$ Gleichungen I und ebenso viele Gleichungen II. Die erste von jeder Art gebe ich als Beispiel an. Die ursprünglichen Gleichungen aus den Beobachtungen vom 20. Mai mit Nadel 1 lauten
\[
\begin{aligned}
i&=67^0\ 11'\ 0''+c-1{,}02880x+1{,}00460y\\
i&=67\ 58\ 46-c-0{,}99473x-1{,}01038y
\end{aligned}
\]
Daraus ergeben sich die abgeleiteten Gleichungen
\[
\begin{aligned}
i&=67^0\ 34'\ 53''-1{,}01176x-0{,}00289y\quad.\ .\ .\ .\ .\ .\quad\text{(I)}\\
c&=+1433''+0{,}01703x-1{,}00749y\quad.\ .\ .\ .\ .\ .\quad\text{(II)}
\end{aligned}
\]
Um die in Artikel 8 angedeutete Prüfung vornehmen zu können, habe ich den Gleichungen I jedoch noch ein Glied hinzugefügt, indem ich $i+e$ statt $i$ schreibe. Dabei bezeichnet $e$ den etwaigen konstanten\HistoricalNote{485} Fehler der Nadel 1. Bei den Nadeln 2, 3 und 4 soll der vermutete konstante Fehler mit $e',e'',e'''$ bezeichnet werden.

Auf diese Weise enthalten die $48$ Gleichungen I insgesamt $36$ unbekannte Größen, nämlich die Inklinationen an den 24 Beobachtungstagen und die $12$ Größen $x,y,e,x',y',e',x''$ und so weiter. Zunächst ist jedoch zu bemerken, dass die Glieder, die $y,y',y'',y'''$ enthalten, alle nur sehr kleine Koeffizienten haben und in der abgekürzten Rechnung ganz fehlen. Der größte dieser $48$ Koeffizienten ist gerade $0.00289$ in der obigen Beispielgleichung. Will
\expandafter\def\csname HistoricalNoteText486\endcsname{Die Beobachtung vom 23. September ist diejenige, die oben in den Artikeln $9$–$22$ als Beispiel gedient hat. Die beiden anderen werden unten in Artikel $30$ angeführt.}
\SourceLeaf{486}man den geringen Einfluss von nur wenigen Sekunden doch berücksichtigen, so muss man die Werte von $y,y',y'',y'''$ zuvor auf anderem Wege abgeleitet haben. Für diesen Zweck reichen aber jedenfalls grobe Näherungswerte aus.

\Article{26}
Für diese Ableitung stehen uns nur die Gleichungen II zur Verfügung. In den $12$ Gleichungen dieser Gruppe, die sich auf jeweils eine Nadel beziehen, steht der Buchstabe $c$ aber für unterschiedliche Werte, da sich sein Wert bei jedem Umstreichen ändern kann. Man sieht daher leicht, dass sich diese Werte von $c$ aus den Gleichungen nicht eliminieren lassen und dass man somit \textit{gezwungen} ist, eine etwas gewagte Hypothese zu Hilfe zu nehmen. Meine Hypothese lautet folgendermaßen. Trotz aller erheblichen Schwankungen von $c$ wird sich bei stets gleichem Streichverfahren ein Mittelwert von $c$ ergeben. Ich nehme daher an, dass dieser Mittelwert für beide Pollagen gleich ist. Freilich ist nur ein sehr unvollkommener Ausgleich zu erwarten, wenn die Nadel nur wenige Male umgestrichen wurde. Der so abgeleitete Wert von $y$ wird daher wenig zuverlässig sein. Dieser Unsicherheit kann man nicht entgehen, wenn man die Werte von $c$ nicht mit einem besonderen Apparat ermittelt (siehe oben Artikel 17). Um dieses Prinzip anzuwenden, trennt man für jede Nadel zunächst die Gleichungen II für $B$ als Nordpol von denen, für die $A$ der Nordpol war. Dann teilt man die beiden Mengen jeweils in so viele Gruppen ein, wie es in der jeweiligen Menge verschiedene magnetische Zustände gab. Aus den Gleichungen jeder Gruppe bildet man den Mittelwert (sofern mehrere Gleichungen in dieselbe Gruppe fallen) und aus diesen Teilmitteln wiederum den Mittelwert. Setzt man die so erhaltenen Mittelwerte gleich, so erhält man die Gleichung, die $y$ bestimmt. Zur Erläuterung gebe ich die \textit{abgekürzte} Rechnung für Nadel 1 an, bei der ich \textit{zu diesem Zweck} außer den obigen $12$ Beobachtungen noch drei weitere\HistoricalNote{486} vom 1. August, 7. August und 23. September verwendet habe. Während des ganzen Zeit\SourceLeaf{487}raums wurde die Nadel neunmal umgestrichen, sodass es zehn verschiedene Zustände gab, fünf für jede Pollage.

\Needspace{16\baselineskip}
\begin{center}\textit{Nadel} 1, $B$ \textit{Nord}\end{center}
\[
\begin{array}{rr|l}
&&\qquad c+y=\\\hline
\text{Mai}&20&.\ .\ .\quad+1433''\\
&\begin{array}{r}24\\31\end{array}&\left.\begin{array}{r}+871''\\+950\end{array}\right\}+910\\
\text{Juni}&5&.\ .\ .\quad+675\\
\begin{array}{r}\text{Juli}\\\phantom{\text{Juli}}\\\text{Aug.}\end{array}&\begin{array}{r}8\\9\\1\end{array}&\left.\begin{array}{r}+739\\+711\\+720\end{array}\right\}+723\\
\begin{array}{r}\phantom{\text{Aug.}}\\\text{Sept.}\end{array}&\begin{array}{r}7\\23\end{array}&\left.\begin{array}{r}+584\\+528\end{array}\right\}+556\\\cline{3-3}
\multicolumn{2}{r}{\text{Mittel}}&c+y=+859''
\end{array}
\]
\Needspace{16\baselineskip}
\begin{center}\textit{Nadel} 1, $A$ \textit{Nord}\end{center}
\[
\begin{array}{rr|l}
&&\\\hline
\text{Mai}&\begin{array}{r}21\\22\end{array}&\left.\begin{array}{r}-653''\\-592\end{array}\right\}-623''\\
\text{Juni}&\begin{array}{r}2\\4\end{array}&\left.\begin{array}{r}-623\\-734\end{array}\right\}-678\\
\text{Juli}&\begin{array}{r}6\\7\end{array}&\left.\begin{array}{r}-730\\-608\end{array}\right\}-669\\
\text{Aug.}&\begin{array}{r}1\\7\end{array}&\left.\begin{array}{r}-720\\-785\end{array}\right\}-752\\
\text{Sept.}&23&.\ .\ .\quad-680\\\cline{3-3}
\multicolumn{2}{r}{\text{Mittel}}&c-y=-680
\end{array}
\]
woraus also $y=+769''$ folgt. Die nicht abgekürzte Rechnung ergab
\[
\begin{aligned}
\text{für }B\text{ Nord},\quad c&=+859+0.00102x-1.00290y\\
\text{für }A\text{ Nord},\quad c&=-680+0.00082x+0.99915y
\end{aligned}
\]
woraus
\[
y=+769''+0.00097x
\]
folgt. Auf dieselbe Weise erhält man für die drei anderen Nadeln
\SourceLeaf{488}\[
\begin{aligned}
y'&=+456''-0.00192x'\\
y''&=-101+0.00134x''\\
y'''&=+1107+0.00224x'''
\end{aligned}
\]
Die Schwankungen der Werte von $c$ erreichen bei Nadel 1 $14\frac12$ Minuten, bei den Nadeln 2 und 3 $4\frac12$ Minuten und bei Nadel 4 $10$ Minuten. Damit man dem Umstand, dass bei Nadel 1 gerade am ersten Beobachtungstag der am stärksten abweichende Wert vorkommt, keine besondere Bedeutung beimisst, möchte ich noch bemerken, dass die Pole sowohl dieser Nadel als auch der übrigen Nadeln schon vor den hier mitgeteilten Beobachtungen oft umgekehrt worden waren, stets mit derselben Sorgfalt und denselben Streichmitteln.

\Article{27}
Nachdem die Werte von $y,y',y'',y'''$ in die Gleichungen I eingesetzt sind, bleiben darin noch $32$ unbekannte Größen. Zieht man dann jeweils die beiden Gleichungen, die für die Beobachtungen desselben Tages gelten, voneinander ab, so erhält man $24$ neue Gleichungen, die nur die acht unbekannten Größen $x,x',x'',x''',e,e',e'',e'''$ enthalten. Die letzten vier kommen aber nur in den Differenzen von je zweien vor. Setzt man daher
\[
\begin{aligned}
e'-e&=d'\\
e''-e&=d''\\
e'''-e&=d'''
\end{aligned}
\]
so bleiben nur sieben unbekannte Größen. Die Koeffizienten von $d',d'',d'''$ sind darin alle $+1$ oder $-1$, und die Koeffizienten von $x,x',x'',x'''$ weichen alle nur sehr wenig von $+1$ oder $-1$ ab. Um die Werte der sieben unbekannten Größen nach der Methode der kleinsten Quadrate zu bestimmen, kann man bei der Bildung der Normalgleichungen für $x,x',x'',x'''$ die Multiplikation mit den jeweiligen Koeffizienten bedenkenlos weglassen. Zur Bildung aller sieben Normalgleichungen ist daher nur einfache Addition nötig. Auf diese Weise haben sich die folgenden Normalgleichungen ergeben
\SourceLeaf{489}\[
\begin{aligned}
0&=+4804''+12.00266x-0.00708x'+0.01900x''\\
0&=-5806+0.01559x+12.01005x'-0.00072x'''\\
0&=-3228+0.00145x+12.00544x''+0.04561x'''\\
0&=-5267+0.01786x'-0.00489x''+12.00343x'''\\
0&=-297+0.02717x+0.11088x'-0.04723x'''\\
 &\qquad-12d'+4d'''\\
0&=-241+0.06326x+0.05839x''-0.08085x'''\\
 &\qquad-12d''+8d'''\\
0&=+254-0.02682x'-0.02676x''+0.12808x'''\\
 &\qquad+4d'+8d''-12d'''
\end{aligned}
\]
und daraus die Werte
\[
\begin{aligned}
x&=-400''\\
x'&=+484\\
x''&=+267\\
x'''&=+438\\
d'&=-22\\
d''&=-23\\
d'''&=+1
\end{aligned}
\]
Anstelle der drei letzten kann man auch
\[
\tfrac14(e+e'+e''+e''')=\varepsilon
\]
setzen und schreiben
\[
\begin{aligned}
e&=+11''+\varepsilon\\
e'&=-11+\varepsilon\\
e''&=-12+\varepsilon\\
e'''&=+12+\varepsilon
\end{aligned}
\]
Der gemeinsame Teil $\varepsilon$ lässt sich aus den verfügbaren Daten offenbar nicht bestimmen. Setzt man die gefundenen Werte von $x,e,x',e'$ und so weiter in die (von $y,y'$ und so weiter bereits befreieten) Gleichungen I ein und lässt $\varepsilon$ weg, so erhält man die folgenden $48$ Inklinationen
\SourceLeaf{490}\Needspace{29\baselineskip}
\begin{center}
\begin{tabular}{rr|c|rrr||c|rrr}
&&Nadel&&&&Nadel&&&\\\hline
Mai&20&$1$&$67^0$&$41'$&$25''$&$2$&$67^0$&$39'$&$12''$\\
&21&&&$39$&$21$&&&$39$&$31$\\
&22&&&$39$&$51$&&&$39$&$22$\\
&24&&&$37$&$43$&&&$40$&$21$\\
&31&&&$40$&$17$&&&$39$&$17$\\
Juni&2&&&$36$&$39$&&&$38$&$16$\\
&4&&&$37$&$31$&&&$37$&$0$\\
&5&&&$41$&$56$&&&$39$&$48$\\
&8&$3$&&$44$&$12$&$4$&&$43$&$14$\\
&9&&&$37$&$27$&&&$38$&$15$\\
&11&&&$36$&$27$&&&$38$&$1$\\
&16&&&$37$&$6$&&&$38$&$17$\\
&18&&&$41$&$12$&&&$40$&$48$\\
&22&&&$38$&$5$&&&$37$&$51$\\
&23&&&$40$&$6$&&&$40$&$2$\\
&25&&&$39$&$45$&&&$39$&$49$\\
Juli&6&$1$&&$40$&$42$&$3$&&$41$&$17$\\
&7&&&$41$&$11$&&&$39$&$49$\\
&8&&&$39$&$5$&&&$41$&$3$\\
&9&&&$39$&$12$&&&$39$&$57$\\
&17&$2$&&$39$&$55$&$4$&&$40$&$7$\\
&18&&&$39$&$56$&&&$39$&$31$\\
&19&&&$39$&$35$&&&$38$&$32$\\
&20&&&$39$&$43$&&&$39$&$1$
\end{tabular}
\end{center}

\Article{28}
Die Unterschiede zwischen den beiden Bestimmungen der Inklination an jedem Tag müssen uns nun ein Maß für die Unsicherheit der Beobachtungen selbst geben. Der größte Unterschied (am 24. Mai) beträgt $2'\ 38''$, und die Summe der Quadrate aller $24$ Unterschiede, in Sekunden als Einheit, ist $124389$. Aus den Grundsätzen der Wahrscheinlichkeitsrechnung lässt sich leicht ableiten, dass wir, wenn wir den Beobachtungen mit jeder der vier Nadeln dieselbe Zuverlässigkeit zuschreiben (und es gibt keinen Grund, von dieser Annahme abzugehen), die mittlere Unsicherheit eines aus den Beobachtungen gewonnenen und unseren Rechnungen zugrunde gelegten Werts von $\tfrac12(f+g)$ oder $\tfrac12(f'+g')$, soweit sich das nach unseren Zahlen beurteilen lässt,
\[
=\BareRoot{\frac{124389}{34}}=60''5
\]
\SourceLeaf{491}ansetzen müssen, sofern es nur um die zufälligen oder regellosen Beobachtungsfehler geht. Das Mittel aus zwei solchen Zahlen, die auf voneinander unabhängigen Beobachtungen beruhen, ist daher mit der mittleren Unzuverlässigkeit
\[
=\BareRoot{\frac{124389}{68}}=42''8
\]
behaftet. Dies kann auch als mittlerer Fehler einer auf die übliche Weise bestimmten Inklination gelten (das heißt mit einer Nadel, aber in \textit{beiden} Pollagen), sofern die kleine zu $\tfrac14(f+g+f'+g')$ hinzukommende Korrektur entweder als ganz unmerklich gilt oder auf einer bereits feststehenden Bestimmung von $u$ oder $y$ beruhen kann (siehe Artikel 21). Es versteht sich von selbst, dass diese Fehlerschätzung zunächst nur für dieses Instrument und für Beobachtungen gilt, die unter ganz ähnlichen Umständen gemacht wurden wie die zugrunde liegenden. Bei weniger als acht Einstellungen in jeder Kombination wäre die Zuverlässigkeit geringer, obwohl ich nicht behaupten möchte, dass der mittlere Fehler des Endergebnisses genau umgekehrt proportional zur Quadratwurzel aus der Zahl der mit Pfannen vervielfältigten Einstellungen ist. Andererseits muss ich darauf hinweisen, dass sich die Lager während der gesamten Dauer der obigen Beobachtungen nicht so vollkommen justieren ließen, wie ich es wünschte und wie es mir später mit dem oben (Artikel 5) erwähnten Apparat tatsächlich gelang. Eine mögliche Vergrößerung der Beobachtungsfehler durch unvollkommene Lagerjustierung ist daher in der obigen Zahl bereits enthalten. An einen Einfluss von konstanter Größe ist dabei umso weniger zu denken, als an den Lagern sehr oft Änderungen vorgenommen wurden. Ich habe daher Grund zu erwarten, dass künftige Beobachtungen mit demselben Instrument eher noch kleinere Fehler zeigen werden.

Eine besondere Untersuchung, deren Einzelheiten ich hier übergehe, hat außerdem ergeben, dass sich die mittlere Unsicherheit der im vorhergehenden Artikel angegebenen $48$ Inklinationen von der mittleren Unsicherheit der Werte $\tfrac12(f+g)$ nur wenig unterscheidet. Den in Artikel 30 zusammenzustellenden Mitteln aus jedem zusammengehörenden Paar ist nahezu das doppelte Gewicht und damit der mittlere Fehler $42''8$ beizulegen.

\SourceLeaf{492}\Article{29}
Ein besonders bemerkenswertes und willkommenes Ergebnis ist, wie klein die für $e,e',e'',e'''$ gefundenen Werte sind, oder genauer zunächst ihre Unterschiede vom Mittel $\varepsilon$. Eine besondere Untersuchung hat für diese Bestimmungen das $\frac{96}{11}$-Fache des Gewichts von $\tfrac12(f+g)$ ergeben, also als mittlere Unsicherheit $=60''5\BareRoot{\frac{11}{96}}=20''5$. Daraus geht hervor, dass selbst das Vorliegen von Unterschieden zwischen $e,e',e'',e'''$ ganz zweifelhaft bleibt. Da es höchst unwahrscheinlich ist, dass vier Nadeln konstante Fehler von fast genau gleicher Größe haben, darf man annehmen, dass sie gar keine oder nur ganz unmerkliche konstante Fehler haben. Es mag daher fast unnötig erscheinen, die Drehbarkeit der Achsen an zweien von ihnen für weitere Versuche zu nutzen.

Für eine der Nadeln, nämlich Nr.\ 4, stützen schon einige frühere Beobachtungen diesen Schluss. An vier Tagen vom 15. bis 19. Mai waren mit den Nadeln 3 und 4 Beobachtungen in ähnlichen Kombinationen wie später vom 8. bis 25. Juni gemacht worden. Der einzige Unterschied war, dass jedes Teilergebnis nur auf vier statt auf acht Einstellungen beruhte. An Nadel 3 waren die Zapfen in derselben Lage wie später, an Nadel 4 dagegen in einer anderen, da nach dem 19. Mai eine Drehung um etwa einen Quadranten vorgenommen wurde. In derselben Schreibweise wie in Artikel 23 lauten die Beobachtungen

\Needspace{8\baselineskip}
\begin{center}\textit{Beobachtungen mit Nadel} 3.

\begin{tabular}{rr|c|rrr|rrr|l}
Mai&15&$B$&$67^0$&$41'$&$26''$&$67^0$&$44'$&$53''$&$6''16166$\\
&17&$B$&&$43$&$52$&&$45$&$52$&$6.20333$\\
&18&$A$&&$33$&$56$&&$39$&$15$&$6.17781$\\
&19&$A$&&$36$&$8$&&$37$&$8$&$6.19566$
\end{tabular}
\end{center}
\Needspace{8\baselineskip}
\begin{center}\textit{Beobachtungen mit Nadel} 4.

\begin{tabular}{rr|c|rrr|rrr|l}
Mai&15&$A$&$67^0$&$14'$&$28''$&$67^0$&$47'$&$49''$&$5''94332$\\
&17&$B$&$68$&$5$&$39$&&$36$&$36$&$5.92034$\\
&18&$B$&&$3$&$30$&&$36$&$13$&$5.94235$\\
&19&$A$&$67$&$3$&$4$&&$59$&$47$&$5.94663$
\end{tabular}
\end{center}

Die Beobachtungen wurden alle vormittags gemacht.

\SourceLeaf{493}
Für Nadel $3$ wurden bei der Rechnung die oben gefundenen Werte von $x'',y'',e''$ verwendet. Für Nadel $4$ mussten dagegen die Werte von $x''',y''',e'''$ so gut wie möglich aus diesen Beobachtungen selbst abgeleitet werden. Dabei ergab sich
\[
\begin{aligned}
y'''&=-1103''\\
x'''&=+556''\\
e'''-\varepsilon&=+24''
\end{aligned}
\]
Die auf so wenige Beobachtungen gestützte Bestimmung von $y'''$ ist allerdings sehr unsicher. Ihr Einfluss auf die Reduktion von $\tfrac12(f+g)$ bleibt jedoch ganz unbedeutend, da der größte Koeffizient von $y'''$ in den Gleichungen I nur $0.00341$ beträgt. Für $i$ ergeben sich dann folgende Werte
\begin{center}
\begin{tabular}{rr||rrr||rrr}
\multicolumn{2}{c||}{}&\multicolumn{3}{c||}{Nadel $3$}&\multicolumn{3}{c}{Nadel $4$}\\\hline
Mai&15&$67^0$&$38'$&$57''$&$67^0$&$40'$&$0''$\\
&17&&$40$&$36$&&$41$&$36$\\
&18&&$41$&$15$&&$40$&$15$\\
&19&&$41$&$19$&&$40$&$16$
\end{tabular}
\end{center}
Das Gewicht der Bestimmung von $e'''-\varepsilon$ ist hier nur doppelt so groß wie das Gewicht von $\tfrac12(f+g)$. Da die Beobachtungen selbst erheblich weniger genau sind als die späteren, spricht der jetzt gefundene Wert offenbar ebenso wenig für das Vorliegen eines konstanten Fehlers wie der aus den späteren Beobachtungen abgeleitete.

Die starke Abweichung der Werte von $x'''$ und $y'''$ von den oben (Artikel 26 und 27) gefundenen Werten zeigt nur, dass der drehbare Teil der Nadel für sich genommen seinen Schwerpunkt nicht in der Zapfenachse hat. Daran ist im Übrigen wenig gelegen.

\Article{30}
Ich stelle nun noch die Endergebnisse für die Inklination aus den sämtli\-chen behandelten Beobachtungen zusammen. Dazu nehme ich auch die Ergebnisse der oben bereits erwähnten Beobachtungen vom 1. und 7. August auf, die mit Nadel 1 genau wie am 23. September gemacht wurden. Die Beobachtungen selbst waren folgende
\SourceLeaf{494}
\begin{center}
\begin{tabular}{r||rrr||rrr}
&\multicolumn{3}{c||}{August $1$}&\multicolumn{3}{c}{August $7$}\\\hline
$f$&$67^0$&$20'$&$12''$&$67^0$&$22'$&$41''$\\
$g$&&$44$&$11$&&$42$&$8$\\
$f'$&&$59$&$53$&$68$&$1$&$56$\\
$g'$&&$35$&$53$&$67$&$35$&$46$
\end{tabular}
\end{center}
\Needspace{20\baselineskip}
\begin{center}\textit{Inklinationsbestimmungen}

\begin{tabular}{rrr|rrr||rr|rrr}
1842&Mai&15&$67^0$&$39'$&$28''$&Juni&18&$67^0$&$41'$&$0''$\\
&&17&&$41$&$6$&&22&&$37$&$58$\\
&&18&&$40$&$45$&&23&&$40$&$4$\\
&&19&&$40$&$47$&&25&&$39$&$47$\\
&&20&&$40$&$18$&Juli&6&&$41$&$0$\\
&&21&&$39$&$26$&&7&&$40$&$30$\\
&&22&&$39$&$36$&&8&&$40$&$4$\\
&&24&&$39$&$2$&&9&&$39$&$34$\\
&&31&&$39$&$47$&&17&&$40$&$1$\\
&Juni&2&&$37$&$27$&&18&&$39$&$44$\\
&&4&&$37$&$15$&&19&&$39$&$4$\\
&&5&&$40$&$52$&&20&&$39$&$22$\\
&&8&&$43$&$43$&Aug.&1&&$39$&$57$\\
&&9&&$37$&$51$&&7&&$40$&$26$\\
&&11&&$37$&$14$&Sept.&23&&$40$&$54$\\
&&16&&$37$&$42$&&&&&
\end{tabular}
\end{center}
Das Mittel aus allen $31$ Bestimmungen, ohne einen Unterschied der Gewichte zu berücksichtigen, beträgt
\[
67^0\;39'\;44''
\]
und kann als für den 21. Juni gültig angesehen werden. Das Mittel aus den $24$ Bestimmungen vom 20. Mai bis 20. Juli allein, dem als mittlerer Zeitpunkt der 19. Juni entspricht, beträgt
\[
67^0\;39'\;31''
\]
\Article{31}
Die Unterschiede der Inklinationen für die einzelnen $31$ Tage von ihrem Mittel setzen sich aus der verbleibenden Wirkung der Beobachtungsfehler und den tatsächlichen Schwankungen der Inklination selbst zusammen. Für die
\SourceLeaf{495}einzelnen Tage lassen sich diese Bestandteile zwar nicht voneinander trennen, doch bei einer Reihe mit so vielen Beobachtungen kann man wohl versuchen, den Mittelwert der tatsächlichen Schwankungen abzuschätzen. Dazu habe ich zunächst die Inklinationen unter der Annahme einer regelmäßigen jährlichen Abnahme von $3$ Minuten auf den 21. Juni reduziert und dann die Quadrate der Abweichungen vom Mittelwert addiert. Dividiert man diese Summe $220184$ durch $30$, so ergibt sich $7339.5$ als Quadrat des mittleren Fehlers, den man in Kauf nähme, wenn man aufs Geratewohl eine jener $31$ Inklinationen als die für die Beobachtungszeit gültige mittlere Inklination ansähe. Will man die unterschiedliche Zuverlässigkeit der drei Beobachtungsgruppen berücksichtigen, so liefern die Grundsätze der Wahrscheinlichkeitsrechnung die folgende Gleichung. Dabei bezeichnet man den mittleren Fehler der ersten vier Beobachtungen mit $m'$, den der letzten drei mit $m''$, den der übrigen $24$ mit $m$ und die mittlere Schwankung der Inklination selbst mit $M$
\[
7339{,}5=\frac{24mm+4m'm'+3m''m''}{31}+MM
\]
Für $mm$ wurde oben der Wert $1829.25$ gefunden, oder zumindest kann man diese Zahl als hinreichende Näherung ansehen. Für die sieben anderen Beobachtungen kann man mangels eines sicheren Maßstabs die Zahl der Einstellungen zugrunde legen, aus denen die Ergebnisse abgeleitet wurden, und somit
\[
m'm'=2mm,\qquad m''m''=\tfrac85mm
\]
setzen. Daraus ergibt sich
\[
MM=7339{,}5-\tfrac{184}{155}.1829{,}25=5168
\]
und $M=71''9$.

\Article{32}
Mit demselben Instrument und am selben Ort hatte ich schon im Vorjahr eine Reihe von Inklinationsbeobachtungen gemacht, von denen ich hier nur die Endergebnisse angebe.
\SourceLeaf{496}
\begin{center}
\begin{tabular}{rrrcrrr}
1841&Sept.&22&$.\quad.\quad.$&$67^0$&$40'$&$20''$\\
&&24&&&$40$&$53$\\
&&27&&&$46$&$41$\\
&Okt.&2&&&$42$&$57$\\
&&7&&&$42$&$14$\\
&&10&&&$42$&$40$\\
&&12&&&$43$&$15$\\
&&20&&&$44$&$2$\\
&&20&&&$42$&$5$\\
&&22&&&$42$&$52$\\\hline
Mittel,&Okt.&8&&$67^0$&$42'$&$48''$
\end{tabular}
\end{center}
Die ersten acht Beobachtungen wurden ähnlich wie die diesjährigen angestellt. An jedem Tag kamen zwei Nadeln (Nr.\ 1 und 2) zum Einsatz, ohne dass die Pole zwischen den Beobachtungen umgekehrt wurden. Die beiden letzten Beobachtungen wurden dagegen auf die übliche Weise gemacht, die zweite vom 20. Oktober mit Nadel 4 und die vom 22. mit Nadel 3. Am 27. September und 10. Oktober wurde nachmittags zwischen $3$ und $5$ Uhr beobachtet, an allen übrigen Tagen vormittags. Jede dieser $10$ Inklinationen beruhte auf $16$ Einstellungen. Aus diesem Grund ist ihnen ein entsprechend geringeres Gewicht zuzuerkennen als den Inklinationen von $1842$, die auf $32$, $64$ beziehungsweise $40$ Einstellungen beruhten.

\Article{33}
Alle bisher angegebenen Inklinationen brauchen noch eine kleine gemeinsame Korrektur für den Einfluss, den die Magnetstäbe der Magnetometer in der Sternwarte und im magnetischen Observatorium am Beobachtungsort ausüben. Um die Ergebnisse von diesem Einfluss zu befreien, muss man durchgehend $5''15$ abziehen (siehe \textit{Resultate} für 1840. II. Artikel 6) [S.\ 433 dieses Bandes]

Die \textit{absolute} Zuverlässigkeit der Inklinationsbestimmungen hängt übrigens weiterhin davon ab, ob die Annahme zutrifft, dass das Instrument selbst keine Teile enthält, die magnetisch auf die Nadel wirken können. Bei dem von mir verwendeten Instrument besteht kein Grund zu einer solchen Befürchtung. Einige Beobachtungen, die ich auf die in Artikel 18 erwähnte Weise mit einer belasteten Nadel anstellte, haben immer nur Abweichun\SourceLeaf{497}gen von ein paar Minuten gezeigt. Sie lassen sich ohne Weiteres durch die unvermeidlichen zufälligen Beobachtungsfehler und die tatsächlichen Anomalien der Inklination selbst erklären. Auch die hinreichend gute Übereinstimmung der Werte, die in Artikel 11 für die dort mit $\alpha$ bezeichnete Größe gefunden wurden, spricht gegen solche Störungen. Zum Erkennen sehr kleiner Einflüsse eignen sich solche Prüfungen freilich nicht. Daher muss ich mir die weitere Prüfung mit gründlicheren Mitteln vorbehalten.

\Article{34}
Zum Schluss stelle ich meine Ergebnisse mit einigen älteren Bestimmungen zusammen.
\[
\begin{array}{r|l}
\begin{array}{rrr}1805&\text{Dez.}&\\1826&\text{Sept.}&\end{array}&
\left.\begin{array}{rrr}69^0&29'&\\68&29&26''\end{array}\right\}\ \text{\textsc{von Humboldt}}\\
\begin{array}{rrr}1837&\text{Juli}&1\\&\text{—}&\text{—}\end{array}&
\left.\begin{array}{rrr}67&47&0\\67&53&30\end{array}\right\}\ \text{\textsc{Forbes}}\\
\begin{array}{rrr}1841&\text{Okt.}&8\\1842&\text{Juni}&21\end{array}&
\begin{array}{rrr}67&42&43\\67&39&39\end{array}
\end{array}
\]
Die ersten beiden Beobachtungen habe ich den \textit{Additions} zum XIII. Band der \textit{Voyage aux régions équinoxiales} entnommen (S.\ 152). Die erste wurde mit einem Inklinometer von \textsc{Lenoir}, die zweite mit einem Instrument von \textsc{Gambey} angestellt. Letztere beruht auf Beobachtungen mit zwei Nadeln. Deren Ergebnisse werden dort mit $68^0\;30'\;7''$ und $68^0\;28'\;15''$ angegeben. Damit stimmt das dort angegebene Mittel nicht überein. Vermutlich ist die Zahl für die zweite Nadel durch einen Druckfehler um $30''$ zu klein angegeben. Den Beobachtungsort von $1805$ kenne ich nicht. $1826$ lag er im freien Feld einige hundert Schritte östlich der Sternwarte.

\textsc{Forbes}' Beobachtungen sind in den \textit{Transactions of the Royal Society of Edinburgh}, Bd.\ XV, Teil 1, S.\ 31 und 32 abgedruckt. Sie wurden im Garten der Sternwarte mit einem \textsc{Robinson}schen Instrument angestellt, das kleiner als das hiesige war, und mit zwei Nadeln von $6$ englischen Zoll Länge. Die zweite Nadel hält der Beobachter selbst für die bessere.

Die von \textsc{Mayer} im März 1814 angestellten und in den \textit{Commentationes recent. Soc. Gotting.}, Bd.\ III, S.\ 36 und 37 ange\SourceLeaf{498}führten Beobachtungen wollte ich hier nicht einreihen, da sie überhaupt kein Vertrauen verdienen. Wie unvollkommen das von \textsc{Mayer} verwendete Instrument war, zeigt die von ihm selbst auf S.\ 35 angegebene Probe. Bei unveränderter Stellung des Instruments ergaben zehn wiederholte Einstellungen Unterschiede von mehr als einem Grad. Seine Ergebnisse für die Inklination selbst, die von zwei verschiedenen Tagen stammen, weichen um einen halben Grad voneinander ab.

Ebenso wenig verdiente meine eigene Beobachtung vom 23. Juni 1832, die in der \textit{Intens. vis magneticae terrestris}, Artikel 27, angeführt ist, hier einen Platz. Das liegt sowohl an der Unvollkommenheit des Instruments als auch am Raum in der Sternwarte. Nicht weit entferntes Eisenwerk musste das Ergebnis erheblich beeinflussen und dabei nachweislich eine Zunahme der Inklination bewirken.

Die angeführten Inklinationen lassen sich sehr gut durch die Annahme einer gleichförmigen jährlichen Abnahme von $3$ Minuten oder genauer $3'\;2''3$ in Einklang bringen, wenn man sich bei \textsc{Forbes}' Beobachtungen an das Ergebnis der zweiten Nadel hält. Es bleiben dann nur Abweichungen, die man ohne Weiteres dem Zusammenwirken der Beobachtungsfehler und der Inklinationsschwankungen zuschreiben kann. Da jedoch nach \textsc{Hansteens} Untersuchungen über die Beobachtungen an anderen europäischen Orten die jährliche Abnahme allmählich langsamer geworden ist, wird man die angegebene Zahl nur als einen ungefähr für $1829$ gültigen Mittelwert betrachten können. Die Bestätigung und genauere Bestimmung der Ungleichförmigkeit wird man erst von künftigen Beobachtungen zu erwarten haben.


\par\endgroup
